\documentclass[11pt]{article}
\usepackage[a4paper,top=25mm,bottom=25mm,left=25mm,right=25mm]{geometry}
\usepackage[T1]{fontenc}
\usepackage[utf8]{inputenc}
\usepackage{lmodern}
\usepackage{microtype}
\usepackage{booktabs,tabularx,array,longtable}
\usepackage{seqsplit}
\usepackage{enumitem}
\usepackage{xcolor}
\usepackage{titlesec}
\usepackage{float}
\usepackage{tikz}
\usetikzlibrary{arrows.meta,positioning,fit,calc,shapes.misc}
\PassOptionsToPackage{hyphens}{url}
\usepackage[hidelinks,breaklinks]{hyperref}
\usepackage[labelfont=bf,font=small]{caption}

\urlstyle{same}
\setlist{nosep,leftmargin=*}
\newcommand{\tphe}{T\mbox{-}PHE}
\newcolumntype{Y}{>{\raggedright\arraybackslash}X}
\newcolumntype{L}[1]{>{\raggedright\arraybackslash\footnotesize}p{#1}}
\sloppy
\setlength{\tabcolsep}{4pt}

\definecolor{ocsi}{RGB}{30,60,90}
\titleformat{\section}{\normalfont\Large\bfseries\color{ocsi}}{\thesection}{0.6em}{}
\titleformat{\subsection}{\normalfont\large\bfseries}{\thesubsection}{0.6em}{}
\titleformat{\subsubsection}{\normalfont\normalsize\bfseries\itshape}{\thesubsubsection}{0.6em}{}

\title{}
\begin{document}

\begin{center}
{\footnotesize\textsc{ARAYA Nikah Social Enterprise Co., Ltd. \,\textbullet\, Bangkok, Thailand}}\par
{\scriptsize Preprint --- not peer reviewed \,\textbullet\, Systematic evidence map (scoping review), reported per PRISMA-ScR}\par\vspace{4mm}
{\LARGE\bfseries\color{ocsi} Evidence for trustworthy premarital health education as primary prevention of intimate-partner harm}\par\vspace{1.5mm}
{\large A systematic evidence map}\par\vspace{4mm}
{\normalsize Yaoharee Lahtee$^{1,2}$}\par
{\small $^{1}$ARAYA Nikah Social Enterprise Co., Ltd., Bangkok, Thailand\\
$^{2}$Independent Researcher, Thailand}\par
{\small ORCID 0009-0005-3861-0626 \,\textbullet\, \texttt{arayawedding@gmail.com}}\par\vspace{2mm}
{\small Draft v0.1 \,\textbullet\, 3 September 2026}\par
\end{center}

\vspace{3mm}
\noindent\textbf{Claim boundary.} This is a systematic evidence map (scoping review) of the
literature searched and admitted while building the evidence library behind the Trustworthy
Premarital Health Education (\tphe) programme theory. It is \emph{not} a systematic review of
effectiveness, has no protocol registration (no PROSPERO record), used single-reviewer
screening with an AI-assisted, human-owned adversarial verification pass rather than independent
duplicate screening, applied no risk-of-bias tool at the individual-study level, and performed
no meta-analysis. Nothing in this document is evidence that \tphe{} works. The claim ceiling is
the accepted TIMA/FIMA 2026 abstract and the \tphe{} full-paper draft (\texttt{paper\slash{}main.pdf} in the public repository); this evidence map is a separate, subordinate document that maps what the
literature search behind that programme theory actually found, admitted, and could not find.

\vspace{4mm}
\noindent\textbf{Abstract}\par\vspace{2pt}
\noindent\begin{minipage}{\textwidth}
\small\raggedright\tolerance=10000\emergencystretch=2em\hbadness=10000
\textbf{Background.} Trustworthy Premarital Health Education (\tphe) is a proposed governance
architecture for premarital education, stated as seven testable propositions (P1--P7) about a
decision-control mechanism, a private channel, cultural intelligibility, a stop rule, referral
quality, the gap between conventional and safeguarding outcomes, and repairable records. Before
any of these propositions is tested, the evidence already assembled while building them needs to
be mapped, not narrated.

\textbf{Methods.} Eighteen category-specific, and to that extent structured, literature searches
were run over PubMed (NCBI E-utilities) plus six supplementary reading-list passes, logged as
18 angle files. A single reviewer screened records; an independent AI-assisted verifier pass
re-fetched each PubMed record and assigned a verdict of KEEP, FIX (admitted with a correction),
DROP, or CONTEXT-ONLY (not PubMed-indexed, no evidential weight). No protocol was registered, no
risk-of-bias tool was applied at study level, and no meta-analysis was performed. All counts
below are computed directly from the logged search, verification, library, and knowledge-graph
files by an accompanying script (\texttt{build\-\_prisma.py}); none is estimated.

\textbf{Results.} The search logged 358 candidate records across 18 angle files (160 PubMed
E-utilities queries), reducing to 277 unique PMIDs after de-duplication among the
PMID-bearing candidates. The verifier processed 358 records: 250 KEEP, 97 FIX, 2 DROP, and 9
CONTEXT-ONLY. The library admits 347 category-level entries (313 unique PMIDs, plus 18 admitted
non-PMID rows corresponding to 17 unique guidance documents without a PMID -- one blank-identifier
collision between two distinct WHO documents was resolved by a deterministic fallback key rather
than a naive count of rows; see Limitations) across 18 categories, 330 unique records in total.
Records linked in the knowledge graph
to at least one of the seven propositions numbered 109 unique records, distributed unevenly:
44 linked to P1, 37 to P5, 23 to P4, 18 to P2, 16 to P3, and only 6 each to P6 and P7. No record
in the library was PubMed-indexed as directly evaluating a premarital education programme, of
any design, in a Thai or Thai Muslim population.

\textbf{Conclusions.} The evidence library supports the plausibility of a premarital governance
layer and repeatedly documents a knowledge-to-behaviour and process-to-outcome gap in adjacent
literatures (screening, relationship education, cultural-competence training), but it contains
no direct effectiveness evidence for \tphe{} or for any comparable programme, and coverage of
Thai/Thai-Muslim samples, premarital-specific trials, validated decision-control instruments, and
the two proposed extension layers (Layer 0, Layer 2) remains sparse to absent. This map is offered
to pre-specify the scope of a future full systematic review with protocol registration, not as a
substitute for one.
\end{minipage}

\vspace{4mm}
\section{Introduction}

\subsection{Why a map, not a narrative}
The \tphe{} programme theory (full paper, \texttt{paper\slash{}main.pdf} in the public repository) states seven testable
propositions and cites an evidence base assembled across 18 literature-search categories. That
evidence base is narrated once, in prose, in the full paper's Discussion and Appendix~E. A
narrative synthesis is not the same document as a map of what was searched, what was found, what
was admitted, what was verified, and what gap each search pass itself logged. This document is
that map: it reports the literature-mapping process as a scoping review in its own right, against
the PRISMA extension for Scoping Reviews (PRISMA-ScR), and treats every count as something to be
computed from the logged files, not recalled from the narrative.

\subsection{Review questions}
The seven propositions that organise the \tphe{} programme theory are used here as the review
questions, unchanged from the full paper:

\begin{itemize}
\item \textbf{P1.} A governance layer improves decision control beyond content-only education.
\item \textbf{P2.} A routine private channel makes materially relevant concerns more visible.
\item \textbf{P3.} Cultural intelligibility helps only when culture cannot override consent or
safety.
\item \textbf{P4.} A stop rule prevents routine joint education from worsening risk when fear or
coercion appears.
\item \textbf{P5.} Referral quality depends on timeliness, accessibility and follow-through, not
offer alone.
\item \textbf{P6.} Conventional programme success can coexist with \tphe{} failure.
\item \textbf{P7.} Repairability is necessary because relevant information emerges over time.
\end{itemize}

For each, Results \S3.3 charts every record the knowledge graph links to that proposition, its
design, its population, its finding as charted, and whether it supports, complicates, or is
relevant only to the proposition (no record in the library is tagged \emph{contradicts} any
proposition; this absence is itself reported, not filled in).

\section{Methods}

\subsection{Eligibility}
Eligible records were PubMed-indexed literature (any design) or named global-guidance documents
(WHO, USPSTF, NICE, Cochrane, MRC) relevant to one of 18 pre-declared search categories spanning
IPV burden and prevention frameworks, premarital/relationship education, screening and referral,
couple-based work and coercive control, decision-autonomy measurement, Muslim/Thailand/Southeast
Asian populations, multicultural clinical safety, preconception medicine, global guidance,
Thailand and the southern border provinces, premarital knowledge and well-being, culture as a
health determinant, digital self-assessment, peer online support, assessor independence and
check-in models, recognition of coercive control, humility/compassion/self-regulation, and
tier-one Islamic medical ethics (library/library.json \texttt{categories}). Sources outside PubMed
and the named global-guidance bodies were admitted only as \emph{context}, carrying no evidential
weight (library/library.json \texttt{context\_sources}); two identified records were dropped
outright as unverifiable (library/library.json \texttt{not\_admitted}).

\subsection{Information sources and search strategy}
Twelve of the 18 categories were searched directly against PubMed using NCBI E-utilities
(\texttt{esearch}/\texttt{efetch}), with every query string logged in that category's
\texttt{search-*.json} file under \texttt{queries\_run}. Six categories (external-culture,
external-dossier, external-draft, external-multicultural, external-thailand, external-wellbeing)
were built from a supplied external reading list rather than a fresh PubMed search pass, and carry
no \texttt{queries\_run} log; this is stated plainly in Table~\ref{tab:queries}, not smoothed over.
The full set of logged query strings, reproduced verbatim, is in Table~\ref{tab:queries}
(Supplementary search strategy).

% Auto-generated by generate_tex_fragments.py -- do not hand-edit.
\begin{longtable}{@{}L{3.4cm}L{11.0cm}@{}}
\caption{Search strategy: the exact PubMed E-utilities query strings logged per angle (queries\_run in each search-*.json file). Six angle files (external-culture, external-dossier, external-draft, external-multicultural, external-thailand, external-wellbeing) were supplied as logged external reading lists rather than a queries\_run search log, and carry no query strings to reproduce here; this is stated, not estimated.}\label{tab:queries}\\
\toprule
\textbf{Search angle} & \textbf{Query string}\\
\midrule
\endfirsthead
\multicolumn{2}{l}{\textit{Table \thetable{} continued}}\\
\toprule
\textbf{Search angle} & \textbf{Query string}\\
\midrule
\endhead
\bottomrule
\endfoot
assessor-\seqsplit{independence}-checkin-models & \texttt{(conflict of interest[Title\slash{}Abstract] OR \seqsplit{independence}[Title\slash{}Abstract]) AND (\seqsplit{safeguarding}[Title\slash{}Abstract] OR "child protection"[Title\slash{}Abstract] OR risk assessment[Title\slash{}Abstract]) AND (assessor[Title\slash{}Abstract] OR \seqsplit{practitioner}[Title\slash{}Abstract] OR provider[Title\slash{}Abstract])}\\
 & \texttt{("intimate partner violence"[MeSH] OR "domestic violence"[MeSH]) AND ("risk assessment"[Title\slash{}Abstract]) AND (DASH[Title\slash{}Abstract] OR ODARA[Title\slash{}Abstract] OR "Danger Assessment"[Title\slash{}Abstract])}\\
 & \texttt{"warm handoff"[Title\slash{}Abstract] AND (referral[Title\slash{}Abstract] OR "care navigation"[Title\slash{}Abstract])}\\
 & \texttt{"brief contact"[Title\slash{}Abstract] AND (self-harm[Title\slash{}Abstract] OR suicide[Title\slash{}Abstract]) AND (\seqsplit{intervention}[Title\slash{}Abstract] OR postcard[Title\slash{}Abstract])}\\
 & \texttt{"single session \seqsplit{intervention}"[Title\slash{}Abstract] AND (prevention[Title\slash{}Abstract] OR mental health[Title\slash{}Abstract])}\\
 & \texttt{postpartum[Title\slash{}Abstract] AND "check-in"[Title\slash{}Abstract] AND (\seqsplit{intervention}[Title\slash{}Abstract] OR screening[Title\slash{}Abstract])}\\
 & \texttt{(fidelity[Title\slash{}Abstract]) AND ("prevention program"[Title\slash{}Abstract] OR "prevention programme"[Title\slash{}Abstract]) AND (incentive[Title\slash{}Abstract] OR completion[Title\slash{}Abstract])}\\
 & \texttt{("dual \seqsplit{relationship}"[Title\slash{}Abstract] OR "conflict of interest"[Title\slash{}Abstract]) AND (clinician[Title\slash{}Abstract] OR therapist[Title\slash{}Abstract]) AND ethics[Title\slash{}Abstract]}\\
 & \texttt{"dual \seqsplit{relationship}"[Title\slash{}Abstract] AND (assessment[Title\slash{}Abstract] OR evaluation[Title\slash{}Abstract]) AND (bias[Title\slash{}Abstract] OR \seqsplit{objectivity}[Title\slash{}Abstract]) -- 0 results}\\
 & \texttt{"\seqsplit{independent} assessor"[Title\slash{}Abstract] OR "\seqsplit{independent} reviewer"[Title\slash{}Abstract] AND (\seqsplit{safeguarding}[Title\slash{}Abstract] OR "risk assessment"[Title\slash{}Abstract]) -- yielded off-topic results, not used}\\
 & \texttt{"who \seqsplit{administers}"[Title\slash{}Abstract] AND ("risk assessment"[Title\slash{}Abstract] OR screening[Title\slash{}Abstract]) AND (intimate partner violence[Title\slash{}Abstract] OR domestic violence[Title\slash{}Abstract]) -- 0 results}\\
 & \texttt{premarital[Title\slash{}Abstract] AND (education[Title\slash{}Abstract] OR counseling[Title\slash{}Abstract]) AND (outcome[Title\slash{}Abstract] OR \seqsplit{effectiveness}[Title\slash{}Abstract]) -- scanned, no new records selected beyond other angle files already in repo}\\
 & \texttt{"conflict of interest"[Title\slash{}Abstract] AND ("child protection"[Title\slash{}Abstract] OR \seqsplit{safeguarding}[Title\slash{}Abstract] OR "mandatory reporting"[Title\slash{}Abstract])}\\
 & \texttt{"structured \seqsplit{professional} judgment"[Title\slash{}Abstract] AND ("intimate partner violence"[Title\slash{}Abstract] OR "domestic violence"[Title\slash{}Abstract])}\\
 & \texttt{"\seqsplit{independent} donor advocate"[Title\slash{}Abstract] OR "\seqsplit{independent} living donor advocate"[Title\slash{}Abstract]}\\
\\[3pt]
couple-work-safety-coercive-control & \texttt{couples therapy intimate partner violence meta-analysis}\\
 & \texttt{conjoint treatment intimate partner violence exclusion criteria}\\
 & \texttt{coercive control assessment \seqsplit{measurement} instrument}\\
 & \texttt{domestic violence screening couples counseling safety}\\
 & \texttt{intimate partner violence premarital counseling}\\
 & \texttt{conjoint couples treatment domestic violence when \seqsplit{appropriate}}\\
 & \texttt{coercive control intimate partner violence \seqsplit{measurement} scale}\\
 & \texttt{\seqsplit{situational} couple violence versus coercive \seqsplit{controlling} violence typology}\\
 & \texttt{couples therapy \seqsplit{contraindicated} domestic violence safety concerns}\\
 & \texttt{intimate partner violence screening tool clinical practice}\\
 & \texttt{batterer \seqsplit{intervention} program couples counseling outcomes}\\
\\[3pt]
decision-autonomy-measures & \texttt{\seqsplit{reproductive} autonomy scale validation}\\
 & \texttt{\seqsplit{reproductive} coercion scale McCauley OR Miller}\\
 & \texttt{decisional autonomy marriage women decision-making power measure}\\
 & \texttt{informed choice measure family planning validation}\\
 & \texttt{coercive control \seqsplit{recognition} instrument scale}\\
 & \texttt{help-seeking intention behavior intimate partner violence measure}\\
 & \texttt{\seqsplit{psychometric} validation autonomy scale Asia OR Muslim}\\
 & \texttt{marital decision making power scale \seqsplit{psychometric}}\\
 & \texttt{"\seqsplit{reproductive} coercion" AND "scale"}\\
 & \texttt{Upadhyay "\seqsplit{reproductive} autonomy"}\\
 & \texttt{McCauley "\seqsplit{reproductive} coercion"}\\
 & \texttt{"decision control" AND marriage}\\
 & \texttt{"women's autonomy" AND Pakistan OR Bangladesh OR Indonesia OR Malaysia}\\
 & \texttt{"coercive control" AND "help-seeking"}\\
 & \texttt{informed choice \seqsplit{contraception} scale \seqsplit{psychometric}}\\
\\[3pt]
digital-self-assessment-ipv & \texttt{I-DECIDE safety decision aid intimate partner violence}\\
 & \texttt{myPlan safety decision aid intimate partner violence}\\
 & \texttt{online safety decision aid domestic violence app randomized}\\
 & \texttt{anonymous self-screening intimate partner violence \seqsplit{acceptability}}\\
 & \texttt{technology based \seqsplit{intervention} intimate partner violence low and middle income countries}\\
 & \texttt{web-based intimate partner violence screening self-\seqsplit{administered}}\\
 & \texttt{premarital screening intimate partner violence risk assessment}\\
 & \texttt{mobile app intimate partner violence safety planning}\\
 & \texttt{online IPV screening Asia \seqsplit{acceptability}}\\
 & \texttt{Hegarty I-DECIDE randomised controlled trial safety}\\
 & \texttt{Glass myPlan intimate partner violence}\\
 & \texttt{Koziol-McLain iSafe online safety decision aid}\\
 & \texttt{Ford-Gilboe iCAN online safety intimate partner violence}\\
 & \texttt{self-assessment tool sexual \seqsplit{reproductive} health premarital}\\
 & \texttt{decision aid \seqsplit{intervention} increase safety behaviour self-efficacy abused women}\\
 & \texttt{computer based screening intimate partner violence primary care}\\
 & \texttt{scored self-assessment app mental health youth anonymous}\\
\\[3pt]
external-culture-citations & \textit{not derivable from the logged files (no queries\_run log; supplied as an external reading list)}\\
\\[3pt]
external-dossier-citations & \textit{not derivable from the logged files (no queries\_run log; supplied as an external reading list)}\\
\\[3pt]
external-draft-citations & \textit{not derivable from the logged files (no queries\_run log; supplied as an external reading list)}\\
\\[3pt]
external-\seqsplit{multicultural}-clinical-safety & \textit{not derivable from the logged files (no queries\_run log; supplied as an external reading list)}\\
\\[3pt]
external-thailand-deepsouth-citations & \textit{not derivable from the logged files (no queries\_run log; supplied as an external reading list)}\\
\\[3pt]
external-wellbeing-citations & \textit{not derivable from the logged files (no queries\_run log; supplied as an external reading list)}\\
\\[3pt]
humility-compassion-self-regulation & \texttt{(\seqsplit{intellectual} humility[tiab]) AND (\seqsplit{relationship}[tiab] OR relational[tiab] OR \seqsplit{forgiveness}[tiab] OR conflict[tiab]) AND (measure[tiab] OR scale[tiab] OR \seqsplit{association}[tiab])}\\
 & \texttt{relational humility[tiab] AND (\seqsplit{forgiveness}[tiab] OR conflict[tiab] OR marital[tiab] OR couple[tiab])}\\
 & \texttt{cultural humility[tiab] AND (health care[tiab] OR healthcare[tiab] OR outcome[tiab] OR patient[tiab])}\\
 & \texttt{cultural humility[tiab] AND (health \seqsplit{disparities}[tiab] OR patient outcomes[tiab] OR systematic review[pt])}\\
 & \texttt{cultural humility[tiab] AND (patient outcomes[tiab] OR clinical outcomes[tiab] OR trust[tiab]) AND (\seqsplit{intervention}[tiab] OR training[tiab])}\\
 & \texttt{self-compassion[tiab] AND (meta-analysis[pt] OR meta-analysis[tiab]) AND (mental health[tiab] OR depression[tiab] OR anxiety[tiab])}\\
 & \texttt{self-compassion[tiab] AND (\seqsplit{relationship} \seqsplit{satisfaction}[tiab] OR \seqsplit{relationship} \seqsplit{functioning}[tiab] OR romantic \seqsplit{relationship}[tiab])}\\
 & \texttt{empathy[tiab] AND (intimate partner violence[tiab] OR IPV[tiab]) AND (\seqsplit{perpetration}[tiab] OR attitude[tiab])}\\
 & \texttt{(self-control[tiab] OR self-regulation[tiab] OR \seqsplit{conscientiousness}[tiab]) AND (intimate partner violence[tiab] OR partner violence \seqsplit{perpetration}[tiab] OR coercive control[tiab]) AND (meta-analysis[pt] OR meta-analysis[tiab])}\\
 & \texttt{(self-control[tiab] OR low self-control[tiab]) AND (intimate partner violence[tiab] OR dating violence[tiab] OR partner violence \seqsplit{perpetration}[tiab]) AND (meta-analysis[pt] OR meta-analysis[tiab] OR systematic review[pt])}\\
 & \texttt{trust[tiab] AND (romantic \seqsplit{relationship}[tiab] OR marital \seqsplit{relationship}[tiab] OR intimate \seqsplit{relationship}[tiab]) AND (\seqsplit{relationship} \seqsplit{satisfaction}[tiab] OR \seqsplit{relationship} quality[tiab])}\\
 & \texttt{\seqsplit{accountability}[tiab] AND (health care ethics[tiab] OR healthcare ethics[tiab] OR patient trust[tiab])}\\
 & \texttt{\seqsplit{religiosity}[tiab] AND (intimate partner violence[tiab] OR domestic violence[tiab]) AND (systematic review[pt] OR systematic review[tiab])}\\
 & \texttt{Islamic psychology[tiab] AND \seqsplit{intervention}[tiab] AND (randomized controlled trial[pt] OR trial[tiab] OR review[pt])}\\
 & \texttt{(humility[tiab] OR values-based[tiab]) AND (couple[tiab] OR marital[tiab] OR premarital[tiab]) AND (\seqsplit{intervention}[tiab] OR program[tiab] OR education[tiab])}\\
 & \texttt{premarital[tiab] AND (education[tiab] OR \seqsplit{intervention}[tiab] OR program[tiab]) AND (randomized[tiab] OR outcome[tiab] OR \seqsplit{effectiveness}[tiab])}\\
 & \texttt{values-based[tiab] AND (couple[tiab] OR premarital[tiab] OR marriage \seqsplit{preparation}[tiab])}\\
\\[3pt]
ipv-screening-response & \texttt{intimate partner violence screening USPSTF \seqsplit{recommendation}}\\
 & \texttt{intimate partner violence screening WHO guideline clinical handbook}\\
 & \texttt{intimate partner violence screening Cochrane systematic review \seqsplit{effectiveness}}\\
 & \texttt{intimate partner violence referral warm referral health care}\\
 & \texttt{domestic violence screening harms unintended \seqsplit{consequences}}\\
 & \texttt{LIVES first-line support intimate partner violence}\\
 & \texttt{World Health \seqsplit{Organization} intimate partner violence guidelines}\\
\\[3pt]
islamic-medical-ethics-tier-one & \texttt{Q1: Islamic bioethics[tiab] AND (principles[tiab] OR framework[tiab]) AND (medical decision[tiab] OR clinical decision[tiab])}\\
 & \texttt{Q2: Islam[tiab] AND informed consent[tiab] AND (autonomy[tiab] OR family[tiab])}\\
 & \texttt{Q3: Muslim[tiab] AND \seqsplit{religiosity}[tiab] AND (mental health[tiab] OR coping[tiab]) AND systematic review[tiab]  [0 hits]}\\
 & \texttt{Q3b: (\seqsplit{spirituality}[tiab] OR \seqsplit{religiosity}[tiab]) AND Muslim[tiab] AND (systematic review[tiab] OR meta-analysis[tiab])}\\
 & \texttt{Q4: Islamic[tiab] AND (\seqsplit{psychotherapy}[tiab] OR \seqsplit{intervention}[tiab]) AND (randomized[tiab] OR meta-analysis[tiab])}\\
 & \texttt{Q5: Islamic[tiab] AND (domestic violence[tiab] OR intimate partner violence[tiab] OR marital[tiab]) AND review[tiab]}\\
 & \texttt{Q6: Muslim physician[tiab] AND (\seqsplit{professionalism}[tiab] OR ethics[tiab] OR trust[tiab])}\\
 & \texttt{Q7: (FIMA[tiab] OR IMANA[tiab] OR "Islamic medical"[tiab]) AND guideline[tiab]  [0 hits]}\\
 & \texttt{Q7b: Islamic[tiab] AND (medical ethics[tiab] OR \seqsplit{professional} ethics[tiab]) AND physician[tiab]}\\
 & \texttt{Q8: maslaha[tiab] OR (Islamic[tiab] AND "harm avoidance"[tiab])}\\
 & \texttt{Q9: shared decision making[tiab] AND Muslim[tiab]}\\
 & \texttt{Q10: Islamic bioethics[tiab] AND \seqsplit{reproductive}[tiab]}\\
 & \texttt{Q11: premarital[tiab] AND (Muslim[tiab] OR Islamic[tiab])}\\
 & \texttt{Q12: Islam[tiab] AND "family \seqsplit{involvement}"[tiab] AND decision making[tiab]}\\
 & \texttt{Q13: Islamic[tiab] AND "BMC Medical Ethics"[Journal] AND (systematic review[tiab] OR review[\seqsplit{Publication} Type])}\\
 & \texttt{Q14: Islam[tiab] AND "Journal of Medical Ethics"[Journal]}\\
 & \texttt{Q15: Islamic[tiab] AND "The American Journal of Bioethics"[Journal]  [0 hits]}\\
\\[3pt]
\seqsplit{manipulation}-literacy-\seqsplit{inoculation} & \texttt{\seqsplit{psychological} \seqsplit{inoculation} \seqsplit{misinformation} \seqsplit{manipulation} techniques randomized}\\
 & \texttt{prebunking \seqsplit{manipulation} techniques video \seqsplit{inoculation}}\\
 & \texttt{coercive control \seqsplit{recognition} scale measure young people}\\
 & \texttt{\seqsplit{gaslighting} \seqsplit{recognition} emotional abuse education}\\
 & \texttt{dating violence prevention program \seqsplit{recognition} attitude school-based randomized}\\
 & \texttt{Safe Dates program dating violence prevention}\\
 & \texttt{systematic review dating violence prevention \seqsplit{adolescents}}\\
 & \texttt{awareness abuse labelling help-seeking intimate partner violence}\\
 & \texttt{health literacy access understand appraise apply \seqsplit{relationship}}\\
 & \texttt{emotional abuse \seqsplit{recognition} Asian Muslim women}\\
 & \texttt{intimate partner violence \seqsplit{recognition} Thailand}\\
 & \texttt{self-check screening tool intimate partner violence anonymous}\\
 & \texttt{Roozenbeek van der Linden fake news game \seqsplit{inoculation}}\\
 & \texttt{health literacy conceptual model access understand appraise apply Sorensen}\\
 & \texttt{warning signs \seqsplit{recognition} dating violence education program \seqsplit{adolescents}}\\
 & \texttt{labeling abuse help-seeking intimate partner violence \seqsplit{qualitative}}\\
 & \texttt{coercive control awareness scale validation}\\
 & \texttt{bystander \seqsplit{intervention} program dating violence \seqsplit{recognition} training}\\
\\[3pt]
muslim-thailand-southeast-asia & \texttt{(intimate partner violence[MeSH] OR IPV) AND (Muslim OR Islam*) AND (Thailand OR Malaysia OR Indonesia OR "Southeast Asia")}\\
 & \texttt{(forced marriage OR fraudulent marriage OR early marriage OR child marriage) AND Muslim}\\
 & \texttt{(premarital OR pre-marital) AND (education OR course OR counseling OR \seqsplit{counselling}) AND (Islam* OR Muslim)}\\
 & \texttt{help-seeking AND intimate partner violence AND Muslim}\\
 & \texttt{One Stop Crisis Center Thailand}\\
 & \texttt{religious leader AND domestic violence AND response}\\
 & \texttt{\seqsplit{unregistered} marriage OR religious marriage AND Southeast Asia}\\
 & \texttt{child marriage AND (Thailand OR Malaysia OR Indonesia)}\\
 & \texttt{kursus pra nikah OR premarital course Malaysia OR Indonesia}\\
 & \texttt{Rohingya AND marriage}\\
 & \texttt{polygamy AND Southeast Asia}\\
 & \texttt{domestic violence Muslim Malaysia}\\
 & \texttt{faith based \seqsplit{intervention} intimate partner violence}\\
 & \texttt{marriage counseling program evaluation Muslim couples}\\
\\[3pt]
peer-online-support-digital-safe-space & \texttt{("intimate partner violence"[tiab]) AND (online[tiab] OR internet[tiab] OR digital[tiab] OR "social media"[tiab] OR forum[tiab]) AND (support[tiab] OR peer[tiab])}\\
 & \texttt{("technology-\seqsplit{facilitated} abuse"[tiab] OR "digital abuse"[tiab] OR "online \seqsplit{surveillance}"[tiab] OR \seqsplit{cyberstalking}[tiab]) AND (partner[tiab] OR \seqsplit{relationship}[tiab]) AND (systematic[tiab] OR review[tiab] OR abuse[tiab])}\\
 & \texttt{("peer support"[tiab]) AND (online[tiab] OR digital[tiab] OR internet[tiab]) AND (adolescent[tiab] OR "young people"[tiab] OR youth[tiab]) AND ("mental health"[tiab]) AND (systematic review[tiab])}\\
 & \texttt{(anonymity[tiab] OR anonymous[tiab]) AND (disclosure[tiab] OR "help-seeking"[tiab] OR "help seeking"[tiab]) AND (online[tiab] OR internet[tiab])}\\
 & \texttt{(moderation[tiab] OR moderator[tiab] OR moderators[tiab]) AND ("online support group"[tiab] OR "online community"[tiab] OR "online \seqsplit{communities}"[tiab] OR "online forum"[tiab] OR "online forums"[tiab]) AND (safety[tiab] OR harm[tiab] OR risk[tiab])}\\
 & \texttt{(adolescent*[tiab] OR "young adult"[tiab] OR "young people"[tiab]) AND (Thailand[tiab] OR Asia[tiab] OR "low- and middle-income"[tiab]) AND ("online community"[tiab] OR "peer support"[tiab] OR "digital health"[tiab])}\\
 & \texttt{(Muslim[tiab] OR Islamic[tiab]) AND (online[tiab] OR digital[tiab]) AND (health[tiab] OR community[tiab])}\\
\\[3pt]
premarital-education-outcomes & \texttt{(premarital OR "marriage \seqsplit{preparation}" OR "marriage and \seqsplit{relationship} education" OR "marital education") AND (counseling OR education OR program) AND (outcome OR \seqsplit{effectiveness})}\\
 & \texttt{"systematic review"[pt] AND (premarital OR "marriage \seqsplit{preparation}" OR "premarital counseling") AND (intimate partner violence OR coercion OR domestic violence) [0 results, \seqsplit{reformulated}]}\\
 & \texttt{systematic review[Title\slash{}Abstract] AND (premarital[Title\slash{}Abstract] OR marriage \seqsplit{preparation}[Title\slash{}Abstract]) AND (violence[Title\slash{}Abstract] OR coercion[Title\slash{}Abstract])}\\
 & \texttt{"Marriage"[Mesh] AND "Counseling"[Mesh] AND ("Intimate Partner Violence"[Mesh] OR "Domestic Violence"[Mesh])}\\
 & \texttt{Hawkins AND "marriage education" AND meta-analysis [0 results, \seqsplit{reformulated}]}\\
 & \texttt{marriage education[Title\slash{}Abstract] AND meta-analysis[Title\slash{}Abstract]}\\
 & \texttt{(faith-based OR religious OR church OR mosque OR clergy) AND ("premarital counseling" OR "marriage \seqsplit{preparation}" OR "marriage education") AND outcomes}\\
 & \texttt{premarital counseling AND help seeking}\\
 & \texttt{premarital education randomized controlled trial}\\
 & \texttt{"Premarital \seqsplit{Examinations}"[Mesh] OR "premarital screening" AND \seqsplit{consanguinity} AND counseling [genetics angle, used only for pool scan, not selected]}\\
\\[3pt]
primary-prevention-frameworks & \texttt{RESPECT+\allowbreak{}women+\allowbreak{}intimate+\allowbreak{}partner+\allowbreak{}violence+\allowbreak{}prevention+\allowbreak{}WHO}\\
 & \texttt{intimate+\allowbreak{}partner+\allowbreak{}violence+\allowbreak{}prevention+\allowbreak{}programme+\allowbreak{}theory+\allowbreak{}logic+\allowbreak{}model}\\
 & \texttt{SASA!+\allowbreak{}community+\allowbreak{}\seqsplit{mobilization}+\allowbreak{}intimate+\allowbreak{}partner+\allowbreak{}violence}\\
 & \texttt{global+\allowbreak{}burden+\allowbreak{}intimate+\allowbreak{}partner+\allowbreak{}violence+\allowbreak{}prevalence+\allowbreak{}WHO+\allowbreak{}2021}\\
 & \texttt{faith+\allowbreak{}based+\allowbreak{}\seqsplit{intervention}+\allowbreak{}intimate+\allowbreak{}partner+\allowbreak{}violence+\allowbreak{}prevention}\\
 & \texttt{\seqsplit{empowerment}+\allowbreak{}\seqsplit{relationship}+\allowbreak{}skills+\allowbreak{}program+\allowbreak{}intimate+\allowbreak{}partner+\allowbreak{}violence}\\
 & \texttt{RESPECT+\allowbreak{}framework+\allowbreak{}preventing+\allowbreak{}violence+\allowbreak{}against+\allowbreak{}women+\allowbreak{}\seqsplit{implementation}}\\
 & \texttt{global+\allowbreak{}estimates+\allowbreak{}prevalence+\allowbreak{}violence+\allowbreak{}against+\allowbreak{}women+\allowbreak{}WHO+\allowbreak{}Sardinha}\\
 & \texttt{global+\allowbreak{}burden+\allowbreak{}intimate+\allowbreak{}partner+\allowbreak{}violence+\allowbreak{}2021+\allowbreak{}Global+\allowbreak{}Burden+\allowbreak{}of+\allowbreak{}Disease}\\
 & \texttt{\seqsplit{Indashyikirwa}+\allowbreak{}couples+\allowbreak{}\seqsplit{intervention}+\allowbreak{}intimate+\allowbreak{}partner+\allowbreak{}violence+\allowbreak{}Rwanda}\\
 & \texttt{Global+\allowbreak{}Burden+\allowbreak{}of+\allowbreak{}Disease+\allowbreak{}2023+\allowbreak{}intimate+\allowbreak{}partner+\allowbreak{}violence+\allowbreak{}women}\\
 & \texttt{GBD+\allowbreak{}2021+\allowbreak{}intimate+\allowbreak{}partner+\allowbreak{}violence+\allowbreak{}disease+\allowbreak{}burden+\allowbreak{}study}\\
\\[3pt]
\end{longtable}


\subsection{Selection of sources of evidence}
A single reviewer ran each search pass, read the abstract of each candidate record, and drafted a
finding and a relevance annotation. An independent verifier pass (AI-assisted, human-owned)
re-fetched every PubMed record by PMID via NCBI E-utilities (no web fetch) and checked title,
year, the finding against the abstract, and whether the relevance annotation overstated the
abstract; WHO/USPSTF/NICE/Cochrane items not indexed in PubMed were checked by title or
identifier search only, not against the live document, and this is flagged per record. Each
record received one verdict: KEEP (admitted as drafted), FIX (admitted with the verifier's
correction, shown verbatim in the record), DROP (not admitted), or CONTEXT-ONLY (not
PubMed-indexed, retained only as context, no evidential weight). This is single-reviewer
screening with an independent adversarial verification pass, not independent duplicate
screening by two reviewers as a full systematic review would require; this is the process's
central limitation and is stated as such throughout, not only here.

\subsection{Data charting items}
For every admitted record, the library charts: identifier (PMID or, for a guidance document, its
own identifier), category, first author, title, year, journal, design, evidence level, population
or setting, finding, a relevance annotation linking the record to the \tphe{} programme theory,
the verifier's verdict and note, and (where applicable) the other categories the same record was
also placed in. The knowledge graph (\texttt{library/kg/graph.json}) additionally links each
record to zero or more of the seven propositions (P1--P7), by a relation of \texttt{supports},
\texttt{complicates}, \texttt{relates\_to} (relevant only), or \texttt{contradicts}; to the
safeguards and outcome domains it measures or discusses; and to its category.

\subsection{Synthesis approach}
No meta-analysis was performed and none is appropriate: the library spans burden estimates,
guidance documents, systematic reviews, single RCTs, qualitative studies, psychometric validation
studies and narrative reviews across 18 categories with no common outcome metric. Synthesis is
narrative, organised two ways: by category (Results \S3.2) and by proposition P1--P7 (Results
\S3.3, one table per proposition, listing every linked record with its design, population, finding
as charted, and its direction relative to the proposition). No risk-of-bias tool was applied at
the individual-study level; a future full systematic review of any one proposition would apply
ROBIS or AMSTAR-2 to included systematic reviews and Cochrane RoB~2 to included randomised trials,
neither of which was in scope for a scoping-level evidence map.

\subsection{Limitations of the process}
(1) Single-reviewer screening throughout, with an AI-assisted independent verifier pass as the
only check, not independent duplicate human screening. (2) No protocol was registered before the
search began; no PROSPERO record exists. (3) Search scope was fixed to PubMed-indexed literature
and five named global-guidance bodies (WHO, USPSTF, NICE, Cochrane, MRC); grey literature, Thai
national databases such as ThaiJO/TCI, and non-English sources were excluded by design and are
logged only where a search pass happened to attempt and fail to resolve one (library/library.json
\texttt{context\_sources}). (4) No risk-of-bias assessment was applied at the individual-study
level. (5) No meta-analysis was performed or attempted. (6) The 18 category searches were run
across different sessions with different query strategies (12 with a logged PubMed query set, 6 as
supplied reading lists), so search sensitivity is not uniform across categories; this is a stated
limitation, not something the counts in \S3.1 attempt to correct for. (7) design and
evidence-level fields were not populated for the six externally-sourced categories
(07--12) in the logged data, so the design/evidence-level distribution in Table~\ref{tab:design}
undercounts those 96 entries as blank rather than reconstructing a design label that was never
charted. (8) \textbf{Data-quality issue in the source library, resolved downstream, not upstream.}
Two genuinely different admitted documents -- a WHO 2013 clinical/policy guideline and a distinct
WHO 2014 clinical handbook, both charted in \texttt{library/library.json} category 03 -- were left
with a blank \texttt{identifier} field, so both would naively collide onto the same
identifier-less key and only one (the 2014 handbook) would ever appear in the rendered tables,
silently dropping the 2013 guideline's own finding from every proposition table it is linked to.
\texttt{library.json} itself was not edited (per this repository's separation between the public
\texttt{library/} and the build scripts released with it); instead, both \texttt{build\_prisma.py}
and \texttt{generate\_tex\_fragments.py} resolve the collision with a deterministic fallback key
(PMID if present, otherwise a slug of first\_author+year+title), so the two documents are counted
and tabulated as the two distinct records they are. This is flagged here as a data-quality issue
of the source library, to be corrected upstream (assigning the 2013 guideline its own non-blank
identifier in \texttt{library.json}) in a future revision of that repository. A related point not
to be conflated with the identifier fix: \texttt{kg/graph.json} records only one
\texttt{relates\_to} edge from the collided source id to each of P2 and P5, so the two documents'
appearance as separate rows in Tables~\ref{tab:p2} and~\ref{tab:p5} rests on applying that single
knowledge-graph edge to both resolved documents, not on two independently coded relations -- the
proposition link itself is exactly as recorded in the knowledge graph, no more and no less.

\section{Results}

\subsection{Selection of sources of evidence (PRISMA flow)}
Figure~\ref{fig:prisma} and Table~\ref{tab:prismaflow} report the PRISMA-ScR flow, computed
directly from the logged files by \texttt{build\_prisma.py}. Records identified through search
(358, sum of hits across 18 angle files, 160 logged PubMed queries) de-duplicate to 277 unique
PMIDs among the PMID-bearing candidates (285 of the 358 candidates carried a PMID at the search
stage; the remaining 73 were guidance documents or non-PMID items in the six externally-sourced
angle files, and are not part of the PMID de-duplication count). All 358 candidate records were
passed to the independent verifier: 250 KEEP, 97 FIX, 2 DROP, 9 CONTEXT-ONLY. The library admits
347 category-level entries, corresponding to 330 unique records (313 unique PMIDs plus 17 admitted
guidance documents without a PMID); the gap between 347 and 330 is records placed in more than one
category (\texttt{also\_in\_category}), each counted once per category it appears in. Of the 18
non-PMID rows, one pair (a WHO 2013 clinical/policy guideline and a distinct WHO 2014 clinical
handbook, both charted with a blank \texttt{identifier} field in \texttt{library.json}) would
naively collide onto the same identifier-less key; both scripts resolve this with a deterministic
fallback key (PMID if present, otherwise a slug of first\_author+year+title) so the two documents
are counted, and tabulated, separately -- see Limitations.

\begin{figure}[H]
\centering
\begin{tikzpicture}[
  node distance=7mm,
  box/.style={draw, rounded corners, align=center, text width=9.3cm, minimum height=9mm, font=\small},
  arr/.style={-{Latex[length=2mm]}}
]
\node[box] (id) {Records identified through search: \textbf{358}\\(18 angle files; 160 logged PubMed queries)};
\node[box, below=of id] (dedup) {Unique PMIDs after de-duplication: \textbf{277}\\(among 285 PMID-bearing candidates; 8 duplicate PMIDs removed)};
\node[box, below=of dedup] (verify) {Records passed to independent verification: \textbf{358}};
\node[box, below=of verify] (verdict) {Verifier verdicts: KEEP \textbf{250} \,\textbullet\, FIX \textbf{97} \,\textbullet\, DROP \textbf{2} \,\textbullet\, CONTEXT-ONLY \textbf{9}};
\node[box, below=of verdict] (admit) {Admitted entries (category-level rows): \textbf{347}\\(330 unique records: 313 unique PMIDs + 17 non-PMID guidance documents)};
\node[box, below=of admit] (link) {Records linked to at least one proposition (P1--P7)\\in the knowledge graph: \textbf{109}};
\draw[arr] (id) -- (dedup);
\draw[arr] (dedup) -- (verify);
\draw[arr] (verify) -- (verdict);
\draw[arr] (verdict) -- (admit);
\draw[arr] (admit) -- (link);
\end{tikzpicture}
\caption{PRISMA-ScR flow. Every number is computed by \texttt{build\_prisma.py} from
library/data/search-*.json, library/data/verify-*.json, library/library.json and
library/kg/graph.json.}
\label{fig:prisma}
\end{figure}

% Auto-generated by generate_tex_fragments.py -- do not hand-edit.
\begin{longtable}{@{}L{13.5cm}r@{}}
\caption{PRISMA-ScR flow counts, computed by build\_prisma.py directly from library/data/search-*.json, library/data/verify-*.json and library/library.json.}\label{tab:prismaflow}\\
\toprule
\textbf{Stage} & \textbf{n}\\
\midrule
\endfirsthead
\multicolumn{2}{l}{\textit{Table \thetable{} continued}}\\
\toprule
\textbf{Stage} & \textbf{n}\\
\midrule
\endhead
\bottomrule
\endfoot
Records identified through structured PubMed E-utilities searches and logged reading-list passes (sum of hits across 18 search-*.json angle files) & 358\\
Total logged PubMed E-utilities queries run (queries\_run arrays) & 160\\
Duplicate PMIDs removed across search files & 8\\
Unique PMIDs identified by search & 277\\
Records passed to the independent verifier & 358\\
Verifier verdict: KEEP & 250\\
Verifier verdict: FIX (correction applied, then admitted) & 97\\
Verifier verdict: DROP (not admitted) & 2\\
Verifier verdict: CONTEXT-ONLY (not PubMed-indexed; no evidential weight) & 9\\
Admitted entries, category-level rows (a record placed in more than one category counts once per category) & 347\\
Admitted unique records (PMID, or identifier for a non-PMID guidance document) & 330\\
Admitted unique PMIDs & 313\\
Admitted entries without a PMID (WHO/NICE/Cochrane/JAKIM/Malaysia-government guidance) & 18\\
Not admitted (DROP, library.json not\_admitted list) & 2\\
Context-only sources logged (not PubMed-indexed, carry no evidential weight) & 9\\
\end{longtable}


\subsection{Characteristics of included records}
Table~\ref{tab:categories} reports admitted entries per category, cross-checked between the count
declared in \texttt{library.json} and a direct recount over the entries themselves (the two agree
for every category). By design/evidence-level keyword count over the free-text fields as logged
(Table~\ref{tab:design}), 41 admitted entries (11.8\% of 347) read as a systematic review, scoping
review, meta-analysis, or Cochrane review by keyword match; 12 (3.5\%) read as a guideline or
recommendation statement, matched only by the two keywords the row label names (``guideline'' and
``recommendation statement''; no broader term such as ``guidance'' is matched, so a record is not
counted here merely for mentioning guidance in passing); and 59 (17.0\%) read as a randomised or
cluster-randomised controlled
trial. These are keyword counts over free text, not a closed design taxonomy applied by a second
coder, and are reported as such. The design and evidence-level fields were not populated for the
96 entries admitted from the six externally-sourced categories (07--12); this is a property of the
logged data, not a search-strategy failure, and is reported as ``not derivable from the logged
files'' rather than reconstructed. Admitted entries span publication years 1991--2026, with 70 of
347 entries (20.2\%) dated 2024 or 2025 and 36 dated 2026 (year-bucketing rule, applied
identically by \texttt{build\_prisma.py}: a \texttt{year} value counts toward a numeric year only
if it is an integer or a clean 4-digit string; every other value -- blank, or an annotated/range
string such as ``2021 (base report); 2023/2024 update per draft claim'' -- is bucketed under
``year not parsed'' and shown as its own row rather than folded into any numeric year; 7 of 347
entries fall in that row. The full year distribution is in the accompanying
\texttt{prisma\_counts.md}).

\begin{table}[H]
\centering\footnotesize
\caption{Admitted entries per category (library.json declared count vs. direct recount over entries; the two agree in every row).}
\label{tab:categories}
\begin{tabularx}{\textwidth}{@{}p{0.6cm}Yr@{}}
\toprule
\textbf{\#} & \textbf{Category} & \textbf{Admitted}\\
\midrule
01 & Burden of intimate-partner violence and primary-prevention frameworks & 20\\
02 & Premarital and relationship education: measured outcomes & 20\\
03 & Screening, first-line response and referral & 21\\
04 & Couple-based work and coercive control & 20\\
05 & Decision autonomy, reproductive coercion and measurement instruments & 18\\
06 & Muslim populations, Thailand and Southeast Asia & 23\\
07 & Multicultural clinical safety and language access & 19\\
08 & Preconception medicine, premarital screening and adjacent guidance & 17\\
09 & Global guidance, frameworks and methods & 29\\
10 & Thailand and the southern border provinces & 15\\
11 & Premarital knowledge and well-being pathways & 7\\
12 & Culture as a health determinant and intercultural coexistence & 9\\
13 & Digital self-assessment and safety decision aids for IPV & 20\\
14 & Peer online support, moderated communities and digital safe spaces & 20\\
15 & Assessor independence, risk assessment and check-in models & 24\\
16 & Recognition of coercive control and psychological manipulation; inoculation & 20\\
17 & Humility, compassion, self-regulation and responsibility & 25\\
18 & Islamic medical ethics and spirituality in Muslim health care (tier-one only) & 20\\
\midrule
\multicolumn{2}{r}{\textbf{Total (category-level rows)}} & \textbf{347}\\
\bottomrule
\end{tabularx}
\end{table}

\begin{table}[H]
\centering\footnotesize
\caption{Design/evidence-level keyword shares over the free-text \texttt{design} and
\texttt{evidence\_level} fields as logged (not a closed taxonomy; a record can match more than one
group, or none). The ``Guideline / recommendation statement'' row matches only the two keywords
its label names -- the bare substrings ``guideline'' and ``recommendation statement'' -- and
deliberately excludes the broader term ``guidance,'' so a record is not counted here merely for
mentioning guidance in passing. The ``RCT / cluster-RCT'' row matches the substrings ``randomised controlled trial,'' ``randomized controlled trial,'' ``rct,'' ``cluster randomi,'' ``randomised trial'' and ``randomized trial'' (case-insensitive).}
\label{tab:design}
\begin{tabularx}{\textwidth}{@{}Yrr@{}}
\toprule
\textbf{Keyword group} & \textbf{n entries} & \textbf{Share of 347 admitted entries}\\
\midrule
Systematic review / meta-analysis / scoping review / Cochrane & 41 & 11.8\%\\
Guideline / recommendation statement & 12 & 3.5\%\\
Randomised or cluster-randomised controlled trial & 59 & 17.0\%\\
\bottomrule
\end{tabularx}
\end{table}

\subsection{Evidence by proposition (P1--P7)}
Tables~\ref{tab:p1}--\ref{tab:p7} chart every record the knowledge graph links to each proposition:
record and design and population as charted in the library, the finding exactly as charted there
(copied verbatim from the record's \texttt{finding} field, or its \texttt{verifier\_note} where
\texttt{finding} was blank), and the direction of that record relative to the proposition
(\texttt{supports}, \texttt{complicates}, or \texttt{relates\_to} shown as ``Relevant only''; no
record in the graph is tagged \texttt{contradicts} any proposition, an absence reported in
\S4.1, not filled in with a record that was not tagged that way). Record counts per proposition are
uneven by construction of the search categories, not by omission: P1 (decision control, the
paper's central mechanism) draws on burden, instrument-validation and community-mobilisation
literatures spanning several categories (44 linked records); P6 (conventional success can coexist
with \tphe{} failure) and P7 (repairable records) are the two propositions argued mainly by
analogy from a small number of directly relevant records (6 each).

\subsubsection*{P1 --- A governance layer improves decision control beyond content-only education}
% Auto-generated by generate_tex_fragments.py -- do not hand-edit.
\begin{longtable}{@{}L{2.0cm}L{2.0cm}L{2.1cm}L{6.2cm}L{1.9cm}@{}}
\caption{P1: A governance layer improves decision control beyond content-only education. Record, design, population/setting, finding as charted (verbatim from the library record), and direction relative to the proposition. n=44 linked records (kg/graph.json edges).}\label{tab:p1}\\
\toprule
\textbf{Record} & \textbf{Design} & \textbf{Population\slash{}setting} & \textbf{Finding (as charted)} & \textbf{Direction}\\
\midrule
\endfirsthead
\multicolumn{5}{l}{\textit{P1 continued}}\\
\toprule
\textbf{Record} & \textbf{Design} & \textbf{Population\slash{}setting} & \textbf{Finding (as charted)} & \textbf{Direction}\\
\midrule
\endhead
\bottomrule
\endfoot
Abramsky T 2014 \cite{ref24} & Pair-matched cluster randomized controlled trial (8 \seqsplit{communities}) & Community members aged 18-49, Kampala, Uganda; baseline n=1,583, follow-up n=2,532 & SASA! was associated with \seqsplit{significantly} lower social acceptance of IPV among women (aRR 0.54, 95\% CI 0.38-0.79) and 52\% lower past-year physical IPV among women (0.48, 95\% CI 0.16-1.39, crossing 1), plus higher likelihood of women receiving supportive community responses to disclosed violence. & Supports\\
\\[3pt]
Alghamdi M 2022 \cite{srp34253076} & \seqsplit{Qualitative} dominant design with \seqsplit{supplemental} \seqsplit{quantitative} data (n=8) & Canadian Muslim immigrant women, 6 countries of origin & Pre-migration adverse childhood \seqsplit{experiences} and family history of IPV, combined with post-migration language\slash{}social-isolation barriers and \seqsplit{internalized} \seqsplit{patriarchal} values, shaped help-seeking and the eventual decision to leave abusive \seqsplit{relationships} among 8 Muslim immigrant women. & Supports\\
\\[3pt]
Blackwood E 2007 \cite{srp17457732} & Historical\slash{}discourse review & Indonesia, state\slash{}religious\slash{}activist discourse, 1980-2000s & Traces how Indonesian state and Islamic discourse shifted from linking normative gender to \seqsplit{heterosexual} marriage toward more direct legal regulation of sexuality via criminal-code proposals, \seqsplit{illustrating} how marriage has functioned as a governance mechanism in national policy discourse. & Supports\\
\\[3pt]
Boyce SC 2023 \cite{srp37316890} & Item response modelling (Partial Credit Model) on population-based dyadic survey & 559 married adolescent-girl-husband dyads (ages 13-18 wives), rural Niger, 2019 & A 5-item, 2-\seqsplit{dimensional} social-norms scale (\seqsplit{challenging} husband authority; general IPV \seqsplit{acceptability}) showed \seqsplit{reliability} and validity; higher '\seqsplit{challenging} husband authority' scores were \seqsplit{statistically} associated with husband \seqsplit{perpetration} of IPV. & Supports\\
\\[3pt]
Böttcher Bettina 2025 \cite{srp41613074} & Cross-cultural validation study & Austrian student population (including men) & The RAS 3-factor structure was supported; decision-making and freedom-from-coercion subscales had \seqsplit{satisfactory} internal \seqsplit{consistency}, while the \seqsplit{communication} subscale showed marginally lower \seqsplit{reliability}; use in men and gender-diverse \seqsplit{individuals} was noted as \seqsplit{constrained}, requiring adaptation. & Supports\\
\\[3pt]
Couture-Carron A 2020 \cite{srp31896308} & \seqsplit{Exploratory} \seqsplit{qualitative} study & South Asian Muslim \seqsplit{communities} (diaspora) & Cultural norms \seqsplit{constructing} dating as shameful were found to intensify fear of parental\slash{}community reaction and strengthen attachment to abusive partners, with direct \seqsplit{implications} for disclosure, help-seeking, and the ability to end abusive \seqsplit{relationships}. & Supports\\
\\[3pt]
GBD 2023 IPV and SVAC \seqsplit{Collaborators} 2026 \cite{ref1} & GBD modelling study (systematic analysis, \seqsplit{spatiotemporal} Gaussian process regression + burden-of-proof risk \seqsplit{methodology}) & 204 countries and \seqsplit{territories}, females aged 15+ (IPV) and males\slash{}females (sexual violence against children), 1990-2023 & In 2023, an estimated 608 million (95\% UI 518-724) females aged 15+ had ever been exposed to IPV, \seqsplit{contributing} 18.5 million (8.74-30.0) DALYs; IPV ranked as the 4th leading risk factor for DALYs among females aged 15-49. Anxiety disorders and major depressive disorder were the leading IPV-attributed health outcomes (5.43 million and 3.96 million DALYs \seqsplit{respectively}). & Supports\\
\\[3pt]
Gausman J 2022 \cite{srp36550423} & \seqsplit{Qualitative} \seqsplit{comparative} study, 8 focus group \seqsplit{discussions} (4 per country) & Syrian refugees (Jordan) and Rohingya refugees (Bangladesh), fathers\slash{}mothers\slash{}\seqsplit{adolescents} & In both displaced Muslim-majority \seqsplit{populations}, \seqsplit{participants} described child marriage norms reinforced across \seqsplit{generations}, with girls themselves describing having little agency in marriage decisions, alongside contextual resistance-to-change dynamics. & Supports\\
\\[3pt]
Gausman Jewel 2023 \cite{srp37922257} & \seqsplit{Psychometric} validation on cross-sectional multi-country survey data & Women in Argentina (n=890), Ghana (n=1,114), and India (n=1,130) & Adapted the Mothers Autonomy in Decision-Making (MADM) scale to \seqsplit{contraceptive} \seqsplit{counselling} contexts, producing the FP-ADM scale, and validated its \seqsplit{psychometric} properties across three culturally distinct LMIC settings including a South Asian one (India). (Verifier correction recorded in the library entry.) & Supports\\
\\[3pt]
Han X 2024 \cite{srp38970860} & Cross-sectional, random-effects Poisson regression (n=201) & Married Salar (Turkic-Muslim minority) women, rural China & Among 201 married Salar Muslim women, 65.6\% reported past-year child neglect, 30.8\% reported past-year IPV, and 37.6\% had \seqsplit{experienced} child marriage; child marriage and IPV \seqsplit{victimization} were associated with downstream child-neglect risk. & Supports\\
\\[3pt]
Hawkins AJ 2012 \cite{ref33} & Meta-analysis (148 evaluation reports) & Mixed university\slash{}laboratory and religious-setting MRE programs, largely US & Moderate-dosage programs (9-20 contact hours) had stronger effects than low-dosage (1-8 hours); a \seqsplit{communication}-skills emphasis \seqsplit{strengthened} \seqsplit{communication} outcomes but not \seqsplit{satisfaction}; no evidence that formal \seqsplit{institutionalization} (manuals, trainer \seqsplit{certification}) improved effects, and no \seqsplit{significant} difference between religious-setting and university\slash{}laboratory delivery. & Supports\\
\\[3pt]
Hazar Seda 2024 \cite{srp38339841} & Cross-cultural adaptation and \seqsplit{psychometric} validation (online survey, EFA\slash{}CFA) & 308 married women of \seqsplit{reproductive} age (15-49), Türkiye, recruited via WhatsApp\slash{}Instagram & The \seqsplit{Reproductive} Autonomy Scale retained its 3-subscale (14-item) structure in Turkish; subscale Cronbach's alpha ranged from 0.64-0.92 (total scale 0.66), with the decision-making subscale showing the lowest \seqsplit{reliability} among subscales tested. (Verifier correction recorded in the library entry.) & Supports\\
\\[3pt]
Hill AL 2019 \cite{srp31306331} & Secondary cross-sectional analysis of baseline survey data & 550 sexually active female high school students, USA (school health center attendees, ages 14-19) & 12\% reported recent \seqsplit{reproductive} coercion and 17\% reported adolescent \seqsplit{relationship} abuse; both were associated with \seqsplit{differences} in care-seeking and sexual health behaviors, measured via validated RC and \seqsplit{relationship}-abuse \seqsplit{instruments}. (Verifier correction recorded in the library entry.) & Supports\\
\\[3pt]
Islam MK 2021 \cite{srp34527969} & Mixed methods: cross-sectional survey (n=96) + \seqsplit{qualitative} interviews (n=18) + provider interviews (n=9) & Rohingya adolescent girls (age 10-19) with experience of child marriage, Cox's Bazar refugee camp, Bangladesh & Mean age at first marriage was 15.7 years; over 80\% had given birth or were pregnant within two years of the survey; key drivers of child marriage included perceived physical\slash{}mental maturity norms, social norms, insecurity, and family honour. & Supports\\
\\[3pt]
Islam MK 2021 \cite{ref5} & Cross-sectional survey, \seqsplit{multivariable} logistic regression (n=486) & Rohingya refugee women, Cox's Bazar settlement, Bangladesh & 61.32\% of surveyed women had \seqsplit{experienced} child marriage; child marriage was associated with \seqsplit{significantly} higher odds of justifying husband-\seqsplit{perpetrated} physical abuse (AOR=2.71, 95\%CI 1.78-4.11) and of \seqsplit{experiencing} IPV in the past 12 months (AOR=1.72, 95\%CI 1.13-2.61). & Supports\\
\\[3pt]
Jacubowski N 2008 \cite{ref112} & \seqsplit{Documentary}\slash{}archival analysis and expert-informant interviews & Married women, Indonesia (Muslim-majority context) & \seqsplit{Traditional} practices including polygamy, early marriage and contract marriage (mut'a) were identified as enhancing married Indonesian women's likelihood of HIV \seqsplit{acquisition} within marriage, \seqsplit{challenging} the assumption that marital status is protective. & Supports\\
\\[3pt]
Kurnaz Kaan 2026 \cite{srp42015126} & Cross-sectional \seqsplit{descriptive}-\seqsplit{correlational} survey with path analysis & 480 married women aged 18-49, Türkiye & Used the \seqsplit{Reproductive} Coercion Scale, Fertility Desire Scale, and Gender Roles Attitude Scale together; \seqsplit{reproductive} coercion and \seqsplit{traditional} gender-role attitudes were linked to fertility desire pathways in married women (path analysis). (Verifier correction recorded in the library entry.) & Supports\\
\\[3pt]
Kuswanto H 2024 \cite{ref3} & Cross-sectional, binary logistic regression (nationally \seqsplit{representative}, n=9,333) & Females aged 15-20, Indonesia (2017 Indonesia \seqsplit{Demographic} and Health Survey) & National prevalence of female child marriage in Indonesia was 12.53\%; health insurance status and sex of household head were not \seqsplit{significant} predictors, while other \seqsplit{socioeconomic}\slash{}education factors were associated with child marriage risk. & Supports\\
\\[3pt]
McCollum 2008 \cite{ref46} & narrative review of outcome research and clinical practice & IPV-affected couples in conjoint treatment settings, US clinical literature & Concludes conjoint IPV treatment can be effective and safe only when embedded in a broader community response, with careful screening of couples for inclusion, \seqsplit{modification} of standard conjoint methods to promote safety, and ongoing safety assessment with \seqsplit{contingency} plans for escalation; a 'one-size-fits-all' approach is explicitly rejected given IPV typologies. & Supports\\
\\[3pt]
Mieth K 2025 \cite{srp40091116} & \seqsplit{Qualitative} study (49 in-depth interviews, 16 focus group \seqsplit{discussions}) & Forcibly Displaced Myanmar Nationals (Rohingya) aged 15-24, Cox's Bazar, Bangladesh (post-October 2016 arrivals) & \seqsplit{Participants} reported perceived increases in child marriage rates post-\seqsplit{displacement}, attributed to protection fears, \seqsplit{socioeconomic} need, lost education\slash{}employment \seqsplit{opportunity}, and perceived loosening of legal-age \seqsplit{enforcement} in the camp context. & Supports\\
\\[3pt]
Riches Eleanor 2023 \cite{srp36657958} & Cross-sectional online survey validation (classical test theory + CFA) & 826 women of \seqsplit{reproductive} age, United Kingdom & The RAS was highly acceptable (>97.7\% response rate) with good internal \seqsplit{consistency} (Cronbach's alpha 0.75) and fair-good test-retest \seqsplit{reliability} (ICC 0.67); construct validity was supported via \seqsplit{confirmatory} factor analysis. & Supports\\
\\[3pt]
Rogers MM 2022 \cite{srp36693203} & \seqsplit{Qualitative} interview study & Family medicine providers\slash{}staff at two US academic outpatient centers & Notes that one in five men report lifetime IPV \seqsplit{perpetration} and two in three male \seqsplit{perpetrators} seek routine health services, yet family medicine physicians feel unprepared to screen male patients for IPV \seqsplit{perpetration} due to lack of knowledge\slash{}training, despite validated screening tools existing. (Verifier correction recorded in the library entry.) & Supports\\
\\[3pt]
Rumble L 2018 \cite{ref4} & Cross-sectional secondary analysis, \seqsplit{multivariate} regression (nationally \seqsplit{representative}, n=6,578) & Women aged 20-24, Indonesia (2012 Indonesian \seqsplit{Demographic} and Health Survey \slash{} Adolescent \seqsplit{Reproductive} Health Survey) & In a nationally \seqsplit{representative} Indonesian sample, \seqsplit{approximately} 17\% of women were married before age 18 and 6\% before age 16; the paper notes research on child marriage in Southeast Asia \seqsplit{specifically} was scarce prior to this nationally \seqsplit{representative} \seqsplit{multivariate} analysis. & Supports\\
\\[3pt]
Sardinha L 2022 \cite{ref2} & Systematic review + Bayesian multilevel modelling of prevalence data & 366 eligible studies, 2 million women, 161 countries\slash{}areas, ever-partnered women aged 15-49 & Globally, 27\% (UI 23-31\%) of ever-partnered women aged 15-49 have \seqsplit{experienced} physical or sexual, or both, IPV in their lifetime, with 13\% (10-16\%) in the past year; violence starts early, affecting 24\% (UI 21-28\%) of women aged 15-19 and 26\% (23-30\%) aged 20-24. (Verifier correction recorded in the library entry.) & Supports\\
\\[3pt]
Sharifnia A 2025 \cite{srp39425490} & Meta-\seqsplit{ethnography} (33 \seqsplit{qualitative} studies \seqsplit{synthesized}), 6 databases searched & Muslim women, global (multi-country \seqsplit{qualitative} evidence) & Synthesis of 33 \seqsplit{qualitative} studies identified four major themes - experience\slash{}impact of abuse, risk factors, help-seeking, and coping - and found religion functions both as a risk factor for abuse and a barrier to help-seeking, and \seqsplit{simultaneously} as a coping resource for some women. & Supports\\
\\[3pt]
Stern E 2020 \cite{srp31686618} & \seqsplit{Longitudinal} \seqsplit{qualitative} study & \seqsplit{Heterosexual} couples \seqsplit{participating} in the \seqsplit{Indashyikirwa} curriculum, Rwanda & The curriculum, designed to build \seqsplit{relationship} skills and reduce gender-\seqsplit{inequitable} beliefs\slash{}behaviours\slash{}norms \seqsplit{underpinning} IPV, was associated with reported shifts in individual beliefs, couple \seqsplit{relationship} dynamics and gender norms among \seqsplit{participating} couples over time. (Verifier correction recorded in the library entry.) & Supports\\
\\[3pt]
Stern E 2021 \cite{srp32896204} & \seqsplit{Comparative} \seqsplit{qualitative} case study (secondary thematic analysis) & \seqsplit{Participants} across four IPV-prevention programmes, What Works to Prevent Violence against Women and Girls Programme (Ghana, Rwanda, South Africa, Tajikistan) & Similar cross-context pathways of change were identified: learning and applying \seqsplit{relationship} skills, \seqsplit{participatory} challenge of harmful gender norms with group rapport, and \seqsplit{integration} of economic-\seqsplit{empowerment} activities to reduce drivers of IPV and build self-confidence. & Supports\\
\\[3pt]
Taft 2024 \cite{ref92} & randomized controlled trial, n=138 couples & US military service members and partners on a military \seqsplit{installation} & Service members randomized to the couples-based Strength at Home Couples (SAH-C) prevention program showed greater effect-size reductions than a comparison Supportive Prevention condition across physical, sexual, \seqsplit{psychological}, and \seqsplit{specifically} coercive-control IPV, plus greater reductions in \seqsplit{suicidality}; partner-reported effects were in the same direction but less robust. & Supports\\
\\[3pt]
Thombs BD 2017 \cite{srp28789659} & \seqsplit{Comparative} \seqsplit{methodological} review of guideline \seqsplit{organizations} (includes IPV screening as one of 22 compared \seqsplit{recommendations}) & Cross-national comparison of screening guideline bodies (Canada, UK, US) & Among 22 \seqsplit{questionnaire}-based screening \seqsplit{recommendations} compared across three major guideline bodies (including IPV screening), \seqsplit{recommendations} diverged \seqsplit{substantially}, and direct RCT evidence for the \seqsplit{effectiveness} of the \seqsplit{questionnaire}-based screening itself (as opposed to downstream treatment) was often absent or indirect. & Supports\\
\\[3pt]
Upadhyay UD 2014 \cite{ref16} & Instrument \seqsplit{development} and validation (cross-sectional survey, factor analysis) & 1,892 women at 13 family planning and 6 abortion facilities, USA & Generated 26 candidate items, retained 14 items across 3 subscales (freedom from coercion, \seqsplit{communication}, decision-making) via factor analysis; the freedom-from-coercion and \seqsplit{communication} subscales were inversely associated with \seqsplit{unprotected} sex in the past 3 months, supporting construct validity. & Supports\\
\\[3pt]
\seqsplit{Ratnaningsih} M 2022 \cite{srp36163286} & Cross-sectional survey, bivariate\slash{}ANOVA\slash{}logistic regression (n=1,000 households) & Households, Bone District (South Sulawesi) and Palu\slash{}Sigi\slash{}Donggala (Central Sulawesi), Indonesia & Household financial security, dowry practices and local legal frameworks were \seqsplit{significantly} associated with community \seqsplit{acceptability} of child marriage across two Indonesian provinces where child-marriage rates remain above the national average. (Verifier correction recorded in the library entry.) & Supports\\
\\[3pt]
Ghuman SJ 2003 \cite{ref93} & \seqsplit{Comparative} study, 15 matched Muslim\slash{}non-Muslim community pairs & Married women, India, Malaysia, the \seqsplit{Philippines} and Thailand & Women's autonomy was not \seqsplit{consistently} lower in Muslim than non-Muslim \seqsplit{communities} across the four Asian countries studied (including Malaysia and Thailand), and the \seqsplit{association} between autonomy measures and child mortality was weak both across and within \seqsplit{communities}. & Complicates\\
\\[3pt]
Hawkins AJ 2008 \cite{ref6} & Meta-analysis (86 reports, 117 studies, >500 effect sizes) & Mostly US\slash{}Western MRE program \seqsplit{participants}; \seqsplit{predominantly} White, middle-class samples & \seqsplit{Relationship}-quality effect sizes for \seqsplit{experimental} studies ranged d=.30-.36 and \seqsplit{communication}-skills effects d=.43-.45; moderate-dosage programs \seqsplit{outperformed} low-dosage; the authors flag that racial\slash{}ethnic and economic diversity was too thin to draw reliable \seqsplit{conclusions} for \seqsplit{disadvantaged} couples, and that outcomes \seqsplit{policymakers} care about (\seqsplit{relationship} stability, aggression) were rarely measured. & Complicates\\
\\[3pt]
Markman HJ 2013 \cite{ref9} & Randomized controlled trial (couples via religious \seqsplit{organizations} randomized to PREP vs. naturally occurring premarital services; n=44\slash{}66\slash{}83) & US couples marrying through religious \seqsplit{organizations}, 8-year follow-up on divorce & No overall difference in divorce rates between PREP (clergy-delivered or \seqsplit{professional}-delivered) and usual premarital services; couples with higher premarital negative-\seqsplit{communication} observed pre-\seqsplit{intervention} were MORE likely to divorce if given PREP, while low-negativity couples benefited; a premarital history of physical aggression showed a similar (marginally \seqsplit{significant}) harmful-moderation pattern. & Complicates\\
\\[3pt]
Miller Elizabeth 2016 \cite{ref110} & Cluster randomized controlled trial (ARCHES) & 4,009 females aged 16-29 across 25 family planning clinics (17 clusters), USA & The universal education\slash{}counseling \seqsplit{intervention} (ARCHES) did NOT \seqsplit{significantly} reduce \seqsplit{reproductive} coercion (ARR 1.50, 95\% CI 0.95-2.35) or partner violence (ARR 1.07, 95\% CI 0.84-1.38) at 1-year follow-up, though it improved secondary outcomes such as \seqsplit{recognition} and self-efficacy. & Complicates\\
\\[3pt]
Phonkacha N 2026 \cite{srp42300599} & Cross-sectional national survey, binary probit regression & General population, national survey, Thailand (Muslim identity modeled as a covariate) & Muslim identity, poverty, and residence in southern\slash{}central Thailand correlated with higher IPV tolerance for both genders; internet access (as distinct from device ownership) \seqsplit{significantly} reduced IPV tolerance, suggesting a role for \seqsplit{information} access in shifting attitudes. & Complicates\\
\\[3pt]
Bhandari TR 2014 \cite{srp26982668} & Scale \seqsplit{construction} and validation (24-item pool, \seqsplit{psychometric} testing, cross-validation across two districts) & Women in Rupandehi and Kapilvastu districts, Nepal (South Asian, \seqsplit{patriarchal} marital-household context) & Produced a 23-item women's autonomy scale with strong \seqsplit{psychometric} properties (Cronbach's alpha 0.84, test-retest r=0.87, content validity ratio 0.8), validated \seqsplit{specifically} for South Asian household\slash{}marital decision contexts, showing good convergent and \seqsplit{discriminant} validity. (Verifier correction recorded in the library entry.) & Relevant only\\
\\[3pt]
Clyde TL 2020 \cite{ref115} & Narrative review\slash{}position paper & General discussion of \seqsplit{contemporary} premarital-education practice, US-centric & Reviews social trends (\seqsplit{individualism}\slash{}commitment \seqsplit{ambivalence}, changing marriage attitudes, premarital \seqsplit{relationship} histories, media \seqsplit{environment}) and argues current premarital \seqsplit{interventions} need updated content and structure to remain relevant to engaged couples today; offers four proposals for \seqsplit{redesigning} premarital \seqsplit{interventions}. & Relevant only\\
\\[3pt]
Kyegombe N 2017 \cite{srp27682273} & \seqsplit{Qualitative} dyadic study (in-depth interviews, framework analysis) & Couples in SASA! \seqsplit{intervention} \seqsplit{communities}, Kampala, Uganda & Engagement with SASA! activities and \seqsplit{relationship}-skills\slash{}\seqsplit{communication} learning was associated with reduced conflict and IPV for many couples and increased trust, respect and joint economic problem-solving, though \seqsplit{experiences} and degree of change varied widely across couples. (Verifier correction recorded in the library entry.) & Relevant only\\
\\[3pt]
Park J 2026 \cite{srp39757299} & Quasi-\seqsplit{experimental} study with inverse \seqsplit{probability} of treatment weighting (weighted N=291 women, 228 men) & Couples in a city-wide (Seoul) premarital education program, South Korea, during COVID-19 & \seqsplit{Videoconferencing} delivery of the Seoul Premarital Education Program (S-PEP) was associated with increased marital readiness for both men and women, and with increased \seqsplit{relationship} confidence and \seqsplit{satisfaction} for women \seqsplit{specifically}. (Verifier correction recorded in the library entry.) & Relevant only\\
\\[3pt]
Portnoy 2020 \cite{srp32791967} & \seqsplit{qualitative} study --- patient interviews (n=10) and provider focus groups (n=29) & US Veterans Health \seqsplit{Administration} patients and providers & Identifies convergent and divergent patient\slash{}provider beliefs about barriers and \seqsplit{facilitators} to screening \seqsplit{specifically} for IPV \seqsplit{perpetration} (not just \seqsplit{victimization}), an area with a much thinner evidence base than \seqsplit{victimization} screening. & Relevant only\\
\\[3pt]
Semahegn A 2017 \cite{srp29162117} & Study protocol for a community-based quasi-\seqsplit{experimental}\slash{}cluster trial & 1,269 women planned across \seqsplit{intervention}, active comparator and control groups, Awi zone, \seqsplit{northwestern} Ethiopia & The protocol describes a planned pilot of \seqsplit{translating} existing evidence and policy into community-level domestic-violence prevention via peer education, community \seqsplit{representative} training, husband \seqsplit{involvement} and \seqsplit{integration} with community health extension programmes; no outcome data are reported (protocol only). & Relevant only\\
\\[3pt]
US Preventive Services Task Force 2004 \cite{srp14996680} & guideline \slash{} \seqsplit{recommendation} statement & US general population, family and intimate partner violence screening in primary care & USPSTF found the evidence \seqsplit{insufficient} (I statement) to recommend for or against routine screening \seqsplit{instruments} to detect domestic violence, updating a 1996 \seqsplit{insufficient}-evidence finding; explicitly states the letter grade reflects the quality and magnitude of evidence available at that time, not an assumption of no benefit. (Verifier correction recorded in the library entry.) & Relevant only\\
\\[3pt]
Upadhyay UD 2021 \cite{ref18} & Instrument \seqsplit{development} and validation (\seqsplit{qualitative} item generation + national survey, EFA, logistic regression) & 1,117 sexually active young people aged 15-24, USA (national sample) & Developed a 23-item, 7-subscale Sexual and \seqsplit{Reproductive} \seqsplit{Empowerment} Scale for \seqsplit{Adolescents}\slash{}Young Adults (including 'choice of partners, marriage, and children' as a named subscale); validation via logistic regression supported construct validity. & Relevant only\\
\\[3pt]
\end{longtable}


\subsubsection*{P2 --- A routine private channel makes materially relevant concerns more visible}
% Auto-generated by generate_tex_fragments.py -- do not hand-edit.
\begin{longtable}{@{}L{2.0cm}L{2.0cm}L{2.1cm}L{6.2cm}L{1.9cm}@{}}
\caption{P2: A routine private channel makes materially relevant concerns more visible. Record, design, population/setting, finding as charted (verbatim from the library record), and direction relative to the proposition. n=18 linked records (kg/graph.json edges).}\label{tab:p2}\\
\toprule
\textbf{Record} & \textbf{Design} & \textbf{Population\slash{}setting} & \textbf{Finding (as charted)} & \textbf{Direction}\\
\midrule
\endfirsthead
\multicolumn{5}{l}{\textit{P2 continued}}\\
\toprule
\textbf{Record} & \textbf{Design} & \textbf{Population\slash{}setting} & \textbf{Finding (as charted)} & \textbf{Direction}\\
\midrule
\endhead
\bottomrule
\endfoot
Abramsky T 2014 \cite{ref24} & Pair-matched cluster randomized controlled trial (8 \seqsplit{communities}) & Community members aged 18-49, Kampala, Uganda; baseline n=1,583, follow-up n=2,532 & SASA! was associated with \seqsplit{significantly} lower social acceptance of IPV among women (aRR 0.54, 95\% CI 0.38-0.79) and 52\% lower past-year physical IPV among women (0.48, 95\% CI 0.16-1.39, crossing 1), plus higher likelihood of women receiving supportive community responses to disclosed violence. & Supports\\
\\[3pt]
Azalee M 2025 \cite{srp40116416} & \seqsplit{Qualitative} in-depth interviews (n=6) & Health care workers who were IPV survivors, tertiary hospital, Malaysia & Six HCW-survivors used varied coping strategies (seeking help, \seqsplit{spirituality}, avoidance, self-harm, staying\slash{}leaving); survivors relied mainly on coworkers and sought formal help only when situations became critical, hindered by \seqsplit{misconceptions} about IPV, privacy concerns and fear of workplace gossip. & Supports\\
\\[3pt]
Chen 2013 \cite{ref94} & clinical practice review & primary\slash{}family care settings & Describes RADAR (Remember to ask routinely, Ask directly in private, Document, Assess safety, Review options) as a structured clinical tool for \seqsplit{identifying} and responding to IPV, \seqsplit{emphasizing} private \seqsplit{questioning}, safety assessment (weapons, children, \seqsplit{perpetrator} risk), and advocacy-based referral requiring trained links to community resources. & Supports\\
\\[3pt]
Decker 2017 \cite{ref95} & quasi-\seqsplit{experimental} single-group pretest-posttest + \seqsplit{qualitative} interviews (survey n=132 baseline, n=68 3-month retention; interviews n=35) & women aged 18+ at two US family planning clinics & A brief, trauma-informed, universal IPV\slash{}\seqsplit{reproductive}-coercion assessment and education \seqsplit{intervention} was \seqsplit{implemented} with 65\% of women receiving at least one \seqsplit{intervention} element; the mixed-methods evaluation describes uptake and provider\slash{}patient experience of the assessment process. & Supports\\
\\[3pt]
Ford-Gilboe 2016 \cite{ref97} & instrument-\seqsplit{development} study using pooled secondary data from 5 Canadian studies, n=6,278 adult women; \seqsplit{international} expert Delphi-style item rating & Canadian adult women, pooled multi-study sample & Developed a brief, validated 15-item short form of the Composite Abuse Scale, addressing the problem that most IPV \seqsplit{measurement} 'privileges physical violence, with poor assessment of other \seqsplit{experiences},' which had led to \seqsplit{underestimating} the scope of IPV in \seqsplit{populations}. (Verifier correction recorded in the library entry.) & Supports\\
\\[3pt]
Hegarty 2026 \cite{ref98} & instrument-\seqsplit{development}\slash{}validation study, n=854 adult women + \seqsplit{qualitative} survivor\slash{}expert feedback & Australian women, community survey sample & Developed and \seqsplit{psychometrically} tested a new self-report Coercive-Composite Abuse Scale (C-CAS) designed to capture severity, frequency and pattern of coercive control \seqsplit{holistically}, explicitly because prior measures showed 'limited \seqsplit{understanding} of these patterns of behaviour.' & Supports\\
\\[3pt]
Kurnaz Kaan 2026 \cite{srp42015126} & Cross-sectional \seqsplit{descriptive}-\seqsplit{correlational} survey with path analysis & 480 married women aged 18-49, Türkiye & Used the \seqsplit{Reproductive} Coercion Scale, Fertility Desire Scale, and Gender Roles Attitude Scale together; \seqsplit{reproductive} coercion and \seqsplit{traditional} gender-role attitudes were linked to fertility desire pathways in married women (path analysis). (Verifier correction recorded in the library entry.) & Supports\\
\\[3pt]
Kyegombe N 2018 \cite{srp30613427} & Cross-sectional survey nested within a cluster randomised trial (SASA! study) & Community members aged 18-49 who had witnessed IPV, Kampala, Uganda, 4 years post-\seqsplit{intervention} & SASA! community members who had witnessed IPV were \seqsplit{substantially} more likely to try to help than control-community \seqsplit{counterparts} (57\% vs 31\%); associated factors differed by trial arm, with SASA!-community bystander action linked to \seqsplit{relationship} duration, employment, talking to neighbours and SASA! exposure. & Supports\\
\\[3pt]
MacMillan 2006 \cite{ref96} & cluster randomized trial, n=2,602 eligible women across 6 clinical sites (emergency, family practice, women's health) & English-speaking women aged 18-64 in Ontario, Canada healthcare settings & Compared face-to-face interview, written self-completed \seqsplit{questionnaire}, and \seqsplit{computerized} screening for IPV to determine the optimal screening method in terms of accuracy, \seqsplit{acceptability} and \seqsplit{completeness}. (Verifier correction recorded in the library entry.) & Supports\\
\\[3pt]
O'Doherty L 2015 \cite{ref38} & Cochrane systematic review and meta-analysis of 13 RCTs\slash{}quasi-RCTs (n=14,959) & Women in antenatal, women's health, emergency department and primary care settings, mostly high-income urban settings & Screening increased clinical \seqsplit{identification} of IPV victims (OR 2.95, 95\% CI 1.79-4.87, n=10,074) but there was no evidence of effect on referrals to support services (OR 2.24, 95\% CI 0.64-7.86, only 2 studies, n=1298), no evidence screening reduced women's re-exposure to violence at 3-18 months, and \seqsplit{insufficient} evidence on health outcomes or harm beyond the immediate post-screening period; authors conclude screening alone is \seqsplit{insufficiently} supported to justify universal \seqsplit{implementation} absent stronger referral\slash{}outcome evidence. & Supports\\
\\[3pt]
Onukwugha FI 2025 \cite{ref42} & Quasi-\seqsplit{experimental} matched-pair cluster-controlled trial (n=1617 clients, 20 \seqsplit{intervention} vs 20 comparison clinics) & Family planning and antenatal care clients, Nigeria & Combined WHO LIVES + ARCHES first-line response reduced relative risk of IPV exposure (beta=0.54, 95\% CI 0.30-0.98) and \seqsplit{reproductive} coercion (beta=0.51, 95\% CI 0.28-0.94) over 9 months among women not already \seqsplit{experiencing} violence, but showed no effect on modern \seqsplit{contraceptive} use; authors note LIVES and ARCHES \seqsplit{individually} have shown 'mixed effects' in prior trials. (Verifier correction recorded in the library entry.) & Supports\\
\\[3pt]
Williamson HC 2014 \cite{ref106} & Cross-sectional survey study (N=2,126 married \seqsplit{individuals}) & US married adults, general population sample & Contrary to the assumption that premarital education reduces later need for counseling, receiving premarital education covaried with an INCREASED likelihood of later couples counseling (a 'gateway' effect), with the \seqsplit{association} stronger among African Americans and those with lower income\slash{}education. & Supports\\
\\[3pt]
Williamson HC 2019 \cite{ref107} & Cross-sectional self-report study (N=231 ethnically diverse newlywed couples in low-income \seqsplit{neighborhoods}) & US low-income newlywed couples & Top perceived barriers to seeking couple therapy were cost and \seqsplit{uncertainty} about where to go for help; direct experience with premarital education (or indirect experience via a social-network member's couple therapy) was associated with reduced barriers and greater likelihood of \seqsplit{subsequently} seeking therapy. & Supports\\
\\[3pt]
\seqsplit{Nsengiyumva} R 2025 \cite{srp41454335} & \seqsplit{Qualitative} \seqsplit{descriptive} study (focus group \seqsplit{discussions}, N=40 engaged couples) & Engaged couples recruited at premarital counseling sites (churches, sector \seqsplit{administration} offices), Rwanda & Engaged couples identified specific unmet \seqsplit{information}\slash{}education needs regarding conception \seqsplit{preparedness}, with premarital counseling sites (including religious venues) identified as a natural point of contact; themes included knowledge gaps, motivators for seeking \seqsplit{information}, and desired \seqsplit{educational} content. (Verifier correction recorded in the library entry.) & Supports\\
\\[3pt]
Evans 2016 \cite{ref51} & \seqsplit{qualitative} think-aloud\slash{}cognitive \seqsplit{interviewing} study with a validated \seqsplit{quantitative} scale (Composite Abuse Scale) & women with experience of domestic violence and abuse & Think-aloud \seqsplit{interviewing} revealed systematic \seqsplit{underreporting} on a widely used, validated DVA scale, \seqsplit{particularly} for coercive control, \seqsplit{threatening} behavior, \seqsplit{restrictions} to freedom, and sexual abuse, driven by item-\seqsplit{interpretation} confusion, fear of answering truthfully, and reluctance to self-identify as abused. & Complicates\\
\\[3pt]
Bhandari TR 2014 \cite{srp26982668} & Scale \seqsplit{construction} and validation (24-item pool, \seqsplit{psychometric} testing, cross-validation across two districts) & Women in Rupandehi and Kapilvastu districts, Nepal (South Asian, \seqsplit{patriarchal} marital-household context) & Produced a 23-item women's autonomy scale with strong \seqsplit{psychometric} properties (Cronbach's alpha 0.84, test-retest r=0.87, content validity ratio 0.8), validated \seqsplit{specifically} for South Asian household\slash{}marital decision contexts, showing good convergent and \seqsplit{discriminant} validity. (Verifier correction recorded in the library entry.) & Relevant only\\
\\[3pt]
World Health \seqsplit{Organization} 2013 (no PMID) & Global clinical and policy guideline & Global health-system guidance for responding to women subjected to IPV or sexual violence & \seqsplit{Establishes} WHO's first-line support principles for clinicians (the basis of the later-named LIVES model: Listen, Inquire about needs, Validate, Enhance safety, Support\slash{}refer), recommends against mandatory universal screening in the absence of adequate response systems, and sets minimum \seqsplit{requirements} (privacy, \seqsplit{confidentiality}, trained staff, referral pathways) before any clinical inquiry about violence should occur. & Relevant only\\
\\[3pt]
World Health \seqsplit{Organization} 2014 (no PMID) & Global clinical practice handbook \seqsplit{operationalizing} the 2013 WHO guidelines & Frontline health workers globally, first-line clinical response to women disclosing IPV or sexual violence & \seqsplit{Operationalizes} the LIVES first-line support protocol (Listen, Inquire, Validate, Enhance safety, Support) into a step-by-step clinical tool for frontline providers, with explicit guidance on \seqsplit{confidentiality}, \seqsplit{documentation}, safety planning, and structured referral rather than passive \seqsplit{information} provision. & Relevant only\\
\\[3pt]
\end{longtable}


\subsubsection*{P3 --- Cultural intelligibility helps only when culture cannot override consent or safety}
% Auto-generated by generate_tex_fragments.py -- do not hand-edit.
\begin{longtable}{@{}L{2.0cm}L{2.0cm}L{2.1cm}L{6.2cm}L{1.9cm}@{}}
\caption{P3: Cultural intelligibility helps only when culture cannot override consent or safety. Record, design, population/setting, finding as charted (verbatim from the library record), and direction relative to the proposition. n=16 linked records (kg/graph.json edges).}\label{tab:p3}\\
\toprule
\textbf{Record} & \textbf{Design} & \textbf{Population\slash{}setting} & \textbf{Finding (as charted)} & \textbf{Direction}\\
\midrule
\endfirsthead
\multicolumn{5}{l}{\textit{P3 continued}}\\
\toprule
\textbf{Record} & \textbf{Design} & \textbf{Population\slash{}setting} & \textbf{Finding (as charted)} & \textbf{Direction}\\
\midrule
\endhead
\bottomrule
\endfoot
Alomair N 2020 \cite{ref20} & Systematic review, narrative synthesis (59 studies, 22 countries) & Muslim women, multi-country (worldwide, including Southeast\slash{}South Asian settings) & Across 59 studies from 22 countries, many Muslim women had poor sexual and \seqsplit{reproductive} health (SRH) knowledge and negative attitudes shaping service access; religious and cultural beliefs, family and community were recurring barriers to \seqsplit{contraception} use and SRH service\slash{}\seqsplit{information} access. & Supports\\
\\[3pt]
Boyce SC 2023 \cite{srp37316890} & Item response modelling (Partial Credit Model) on population-based dyadic survey & 559 married adolescent-girl-husband dyads (ages 13-18 wives), rural Niger, 2019 & A 5-item, 2-\seqsplit{dimensional} social-norms scale (\seqsplit{challenging} husband authority; general IPV \seqsplit{acceptability}) showed \seqsplit{reliability} and validity; higher '\seqsplit{challenging} husband authority' scores were \seqsplit{statistically} associated with husband \seqsplit{perpetration} of IPV. & Supports\\
\\[3pt]
Fauk NK 2021 \cite{ref103} & \seqsplit{Qualitative} in-depth interviews (n=92: 52 women, 40 men) & People living with HIV, Belu and Yogyakarta, Indonesia (Islamic and Javanese cultural context) & Cultural practices (bride wealth, spousal-dispute management) and Islamic religious beliefs about husband-wife \seqsplit{relationships}, forbidden premarital sex, and condom use were reported by \seqsplit{participants} as shaping spousal sexual behaviours that \seqsplit{contributed} to HIV \seqsplit{transmission} within marriage. & Supports\\
\\[3pt]
James LE 2021 \cite{srp34770188} & Co-designed, two-phase community-based \seqsplit{intervention} (workshops); RCT-labelled but abstract available describes \seqsplit{intervention} \seqsplit{development}\slash{}testing & Rohingya (Muslim minority, displaced) in Malaysia and Syrian refugees (Muslim-majority) in Lebanon & A social-norms-based, mental-health-integrated community \seqsplit{intervention} was co-designed with Rohingya and Syrian refugee \seqsplit{communities} \seqsplit{specifically} to reduce IPV and encourage help-seeking, using culturally embedded workshop delivery rather than clinical\slash{}individual counseling alone. & Supports\\
\\[3pt]
Meiksin R 2020 \cite{srp32460786} & Trial protocol & Women\slash{}girls seeking \seqsplit{contraceptive} services, community-based clinics, Nairobi, Kenya & Describes ARCHES as a single-session, clinic-based enhanced counseling model that creates \seqsplit{opportunity} for patient disclosure of \seqsplit{reproductive} coercion\slash{}IPV followed by warm referral to local services; explicitly notes that to date, no clinic-based GBV \seqsplit{intervention} had been shown to reduce \seqsplit{reproductive} coercion or unintended pregnancy in LMIC settings, motivating this trial. (Verifier correction recorded in the library entry.) & Supports\\
\\[3pt]
Owen JJ 2011 \cite{srp21355646} & RCT sub-analysis (31 leaders training 118 couples\slash{}236 attendees in a PREP trial) & US couples in a randomized PREP trial, multiple \seqsplit{facilitators} & \seqsplit{Relationship} outcomes varied \seqsplit{significantly} by which leader delivered the training; \seqsplit{facilitators}' aggregated working-alliance quality (as rated by attendees) predicted change in observed \seqsplit{communication} and \seqsplit{relationship} \seqsplit{satisfaction}, over and above curriculum content itself. (Verifier correction recorded in the library entry.) & Supports\\
\\[3pt]
Sarac A 2026 \cite{ref100} & Pilot randomized controlled trial, mixed methods (N=8 couples primary phase + 2 couples \seqsplit{feasibility} phase) & Young premarital couples, Turkey (culturally adapted European program) & A culturally adapted, 6-session Couple Learning Program produced \seqsplit{significant} \seqsplit{improvements} in conflict-resolution strategies (exit, voice, neglect dimensions) with \seqsplit{relationship} \seqsplit{satisfaction} improving only at follow-up (not \seqsplit{immediately} post-\seqsplit{intervention}); \seqsplit{participants} rated the culturally adapted program as relevant and acceptable. & Supports\\
\\[3pt]
Sharifnia A 2025 \cite{srp39425490} & Meta-\seqsplit{ethnography} (33 \seqsplit{qualitative} studies \seqsplit{synthesized}), 6 databases searched & Muslim women, global (multi-country \seqsplit{qualitative} evidence) & Synthesis of 33 \seqsplit{qualitative} studies identified four major themes - experience\slash{}impact of abuse, risk factors, help-seeking, and coping - and found religion functions both as a risk factor for abuse and a barrier to help-seeking, and \seqsplit{simultaneously} as a coping resource for some women. & Supports\\
\\[3pt]
Stern E 2021 \cite{srp32896204} & \seqsplit{Comparative} \seqsplit{qualitative} case study (secondary thematic analysis) & \seqsplit{Participants} across four IPV-prevention programmes, What Works to Prevent Violence against Women and Girls Programme (Ghana, Rwanda, South Africa, Tajikistan) & Similar cross-context pathways of change were identified: learning and applying \seqsplit{relationship} skills, \seqsplit{participatory} challenge of harmful gender norms with group rapport, and \seqsplit{integration} of economic-\seqsplit{empowerment} activities to reduce drivers of IPV and build self-confidence. & Supports\\
\\[3pt]
Tiwari 2015 \cite{srp24860075} & mixed-methods evaluation study, n=200 Chinese women survivors, in-depth interviews + ROC curve analysis & Chinese women survivors of intimate partner violence & \seqsplit{Qualitative} accounts of 91\slash{}200 women were consistent with Dutton and Goodman's coercive-control \seqsplit{conceptualization}; the Chinese Revised \seqsplit{Controlling} Behaviors Scale showed acceptable \seqsplit{discriminative} accuracy (reported via ROC\slash{}AUC) for \seqsplit{identifying} high-control \seqsplit{relationships} in a non-Western cultural sample. & Supports\\
\\[3pt]
Welland C 2010 \cite{ref101} & \seqsplit{Qualitative} study (\seqsplit{demographic} survey + \seqsplit{qualitative} interviews) informing program design & Latino immigrant men who are intimate-partner-violence \seqsplit{perpetrators}, Southern California & A culturally specific, Spanish-language program \seqsplit{integrating} gender roles, \seqsplit{immigration}\slash{}\seqsplit{discrimination} context, marital sexual abuse, and \seqsplit{spirituality} was seen by \seqsplit{participants} as more relevant; the study explicitly frames cultural \seqsplit{sensitivity} as compatible with, not a substitute for, \seqsplit{perpetrator} \seqsplit{accountability}. & Supports\\
\\[3pt]
\seqsplit{Ratnaningsih} M 2022 \cite{srp36163286} & Cross-sectional survey, bivariate\slash{}ANOVA\slash{}logistic regression (n=1,000 households) & Households, Bone District (South Sulawesi) and Palu\slash{}Sigi\slash{}Donggala (Central Sulawesi), Indonesia & Household financial security, dowry practices and local legal frameworks were \seqsplit{significantly} associated with community \seqsplit{acceptability} of child marriage across two Indonesian provinces where child-marriage rates remain above the national average. (Verifier correction recorded in the library entry.) & Supports\\
\\[3pt]
Dagher Myriam 2024 \cite{ref19} & 4-step cultural adaptation (cognitive interviews, cross-sectional \seqsplit{administration} n=339, \seqsplit{confirmatory} factor analysis) & Refugee and non-refugee adolescent girls aged 11-17, Lebanon (Arabic-speaking, majority-Muslim region), including \seqsplit{participants} in an early-marriage \seqsplit{intervention} & A Western-developed sexual and \seqsplit{reproductive} \seqsplit{empowerment} scale was adapted to Arabic; the original 13-item, 4-domain model showed poor fit in \seqsplit{confirmatory} factor analysis, and even after \seqsplit{iteratively} removing 2 items, model fit remained sub-optimal for most domains. & Complicates\\
\\[3pt]
Sikoki B 2024 \cite{ref102} & \seqsplit{Qualitative} (focus groups and individual interviews, feminist analytic lens), fieldwork 2016 & Javanese Muslim \seqsplit{adolescents} and adults, Yogyakarta, Indonesia & School-based SRH education in the Javanese Muslim community was \seqsplit{constrained} by taboo and moral-panic framing around premarital sex; delivery gaps left \seqsplit{adolescents} without reliable sexuality \seqsplit{information} despite peer-counsellor programmes intended to fill this role. & Complicates\\
\\[3pt]
Wood SN 2020 \cite{srp33336831} & Mixed-methods (\seqsplit{exploratory} factor analysis, n=327; in-depth interviews, n=30) & Recent IPV survivors, informal \seqsplit{settlements}, Nairobi, Kenya & The US-developed \seqsplit{Reproductive} Coercion Scale \seqsplit{transferred} reasonably well to the Kenyan context, retaining a 2-factor structure (pregnancy coercion, condom \seqsplit{manipulation}; alpha=0.86); 82\% of IPV survivors reported some \seqsplit{reproductive} coercion. & Complicates\\
\\[3pt]
Muluh N 2023 \cite{ref27} & Two-arm, mixed-methods cluster randomized controlled trial (18-month faith-based, norms-shifting programme) & Women aged 18-35 (n=350) and male partners (n=281), from 10 Muslim and 10 Christian \seqsplit{congregations}, Plateau State, Nigeria & The \seqsplit{Masculinity}, Faith, and Peace (MFP) programme, delivered through faith \seqsplit{congregations}, assessed shifts in individual-, couple- and community-level social norms and IPV incidence in a mixed Muslim\slash{}Christian \seqsplit{congregational} sample; the trial was explicitly designed and powered to evaluate norms-shifting mechanisms rather than knowledge\slash{}attendance alone. (Verifier correction recorded in the library entry.) & Relevant only\\
\\[3pt]
\end{longtable}


\subsubsection*{P4 --- A stop rule prevents routine joint education from worsening risk when fear or coercion appears}
% Auto-generated by generate_tex_fragments.py -- do not hand-edit.
\begin{longtable}{@{}L{2.0cm}L{2.0cm}L{2.1cm}L{6.2cm}L{1.9cm}@{}}
\caption{P4: A stop rule prevents routine joint education from worsening risk when fear or coercion appears. Record, design, population/setting, finding as charted (verbatim from the library record), and direction relative to the proposition. n=23 linked records (kg/graph.json edges).}\label{tab:p4}\\
\toprule
\textbf{Record} & \textbf{Design} & \textbf{Population\slash{}setting} & \textbf{Finding (as charted)} & \textbf{Direction}\\
\midrule
\endfirsthead
\multicolumn{5}{l}{\textit{P4 continued}}\\
\toprule
\textbf{Record} & \textbf{Design} & \textbf{Population\slash{}setting} & \textbf{Finding (as charted)} & \textbf{Direction}\\
\midrule
\endhead
\bottomrule
\endfoot
Couture-Carron A 2020 \cite{srp31896308} & \seqsplit{Exploratory} \seqsplit{qualitative} study & South Asian Muslim \seqsplit{communities} (diaspora) & Cultural norms \seqsplit{constructing} dating as shameful were found to intensify fear of parental\slash{}community reaction and strengthen attachment to abusive partners, with direct \seqsplit{implications} for disclosure, help-seeking, and the ability to end abusive \seqsplit{relationships}. & Supports\\
\\[3pt]
Han X 2024 \cite{srp38970860} & Cross-sectional, random-effects Poisson regression (n=201) & Married Salar (Turkic-Muslim minority) women, rural China & Among 201 married Salar Muslim women, 65.6\% reported past-year child neglect, 30.8\% reported past-year IPV, and 37.6\% had \seqsplit{experienced} child marriage; child marriage and IPV \seqsplit{victimization} were associated with downstream child-neglect risk. & Supports\\
\\[3pt]
Hegarty K 2025 \cite{srp41330529} & \seqsplit{Qualitative} directed content analysis of open-ended survey responses (n=682 women with lived experience of coercive control) & Australian women who had \seqsplit{experienced} coercive control & Findings align with the WHO LIVES framework (Listen-Inquire-Validate-Enhance safety-Support) and the CARE framework (Choice and control-Action and advocacy-\seqsplit{Recognition} and \seqsplit{understanding}-Emotional connection); women wanted specific validating, non-judgmental scripted language from \seqsplit{practitioners} and messaging that \seqsplit{counteracts} coercive-\seqsplit{controlling} tactics rather than generic \seqsplit{acknowledgement}. (Verifier correction recorded in the library entry.) & Supports\\
\\[3pt]
Islam MK 2021 \cite{srp34527969} & Mixed methods: cross-sectional survey (n=96) + \seqsplit{qualitative} interviews (n=18) + provider interviews (n=9) & Rohingya adolescent girls (age 10-19) with experience of child marriage, Cox's Bazar refugee camp, Bangladesh & Mean age at first marriage was 15.7 years; over 80\% had given birth or were pregnant within two years of the survey; key drivers of child marriage included perceived physical\slash{}mental maturity norms, social norms, insecurity, and family honour. & Supports\\
\\[3pt]
James LE 2021 \cite{srp34770188} & Co-designed, two-phase community-based \seqsplit{intervention} (workshops); RCT-labelled but abstract available describes \seqsplit{intervention} \seqsplit{development}\slash{}testing & Rohingya (Muslim minority, displaced) in Malaysia and Syrian refugees (Muslim-majority) in Lebanon & A social-norms-based, mental-health-integrated community \seqsplit{intervention} was co-designed with Rohingya and Syrian refugee \seqsplit{communities} \seqsplit{specifically} to reduce IPV and encourage help-seeking, using culturally embedded workshop delivery rather than clinical\slash{}individual counseling alone. & Supports\\
\\[3pt]
Kalichman SC 2020 \cite{srp32869478} & Commentary \slash{} research-agenda-setting article & Adolescent girls and young women (15-24) in LMICs receiving HIV index-testing\slash{}partner-\seqsplit{notification} services & Highlights that partner-\seqsplit{notification} and index-testing approaches proven safe\slash{}effective for adult men and cohabiting\slash{}married adults have not been adequately studied for younger women with less \seqsplit{relationship} power, who face higher risk of IPV, stigma, and economic coercion when disclosure is required or \seqsplit{facilitated} by a health service. (Verifier correction recorded in the library entry.) & Supports\\
\\[3pt]
Karakurt 2016 \cite{ref45} & systematic review and meta-analysis (6 studies of 1,733 screened citations) & couples in treatment for physical\slash{}\seqsplit{situational} IPV, mostly US clinical\slash{}research samples & After screening 1,733 citations, only 6 studies met inclusion criteria for meta-analysis of couples therapy as a treatment for IPV; \seqsplit{preliminary} pooled data suggest couples therapy can be a viable treatment 'in select situations' but the evidence base is small and \seqsplit{professionals} remain cautious about risk of violent \seqsplit{retaliation} between partners. & Supports\\
\\[3pt]
Kuswanto H 2024 \cite{ref3} & Cross-sectional, binary logistic regression (nationally \seqsplit{representative}, n=9,333) & Females aged 15-20, Indonesia (2017 Indonesia \seqsplit{Demographic} and Health Survey) & National prevalence of female child marriage in Indonesia was 12.53\%; health insurance status and sex of household head were not \seqsplit{significant} predictors, while other \seqsplit{socioeconomic}\slash{}education factors were associated with child marriage risk. & Supports\\
\\[3pt]
McCauley HL 2017 \cite{ref17} & Secondary \seqsplit{psychometric} analysis of pooled \seqsplit{multicenter}\slash{}RCT survey data & 4,674 women aged 16-29 in 24 \seqsplit{Pennsylvania} and 5 California family planning clinics, USA & Confirmed a 2-domain structure (pregnancy coercion, condom \seqsplit{manipulation}) and produced a validated 5-item short-form \seqsplit{Reproductive} Coercion Scale (RCS); recent \seqsplit{reproductive} coercion prevalence was 6.3-6.7\% depending on the scale version used. & Supports\\
\\[3pt]
McCollum 2008 \cite{ref46} & narrative review of outcome research and clinical practice & IPV-affected couples in conjoint treatment settings, US clinical literature & Concludes conjoint IPV treatment can be effective and safe only when embedded in a broader community response, with careful screening of couples for inclusion, \seqsplit{modification} of standard conjoint methods to promote safety, and ongoing safety assessment with \seqsplit{contingency} plans for escalation; a 'one-size-fits-all' approach is explicitly rejected given IPV typologies. & Supports\\
\\[3pt]
Mieth K 2025 \cite{srp40091116} & \seqsplit{Qualitative} study (49 in-depth interviews, 16 focus group \seqsplit{discussions}) & Forcibly Displaced Myanmar Nationals (Rohingya) aged 15-24, Cox's Bazar, Bangladesh (post-October 2016 arrivals) & \seqsplit{Participants} reported perceived increases in child marriage rates post-\seqsplit{displacement}, attributed to protection fears, \seqsplit{socioeconomic} need, lost education\slash{}employment \seqsplit{opportunity}, and perceived loosening of legal-age \seqsplit{enforcement} in the camp context. & Supports\\
\\[3pt]
Nelson HD 2012 \cite{srp22565034} & systematic review (evidence base for USPSTF 2013 update) & Women in health care settings, US-relevant evidence base & Of 15 fair\slash{}good-quality studies evaluating 13 screening \seqsplit{instruments}, 6 were highly accurate; 14 studies found minimal adverse effects from screening, though some women reported discomfort, loss of privacy, emotional distress, and fear of further abuse; counseling \seqsplit{interventions} reduced IPV\slash{}improved outcomes in some \seqsplit{subpopulations} (pregnant women, new mothers, family-planning clients) but a large trial found no \seqsplit{significant} screening-vs-usual-care difference overall. & Supports\\
\\[3pt]
Onukwugha FI 2025 \cite{ref42} & Quasi-\seqsplit{experimental} matched-pair cluster-controlled trial (n=1617 clients, 20 \seqsplit{intervention} vs 20 comparison clinics) & Family planning and antenatal care clients, Nigeria & Combined WHO LIVES + ARCHES first-line response reduced relative risk of IPV exposure (beta=0.54, 95\% CI 0.30-0.98) and \seqsplit{reproductive} coercion (beta=0.51, 95\% CI 0.28-0.94) over 9 months among women not already \seqsplit{experiencing} violence, but showed no effect on modern \seqsplit{contraceptive} use; authors note LIVES and ARCHES \seqsplit{individually} have shown 'mixed effects' in prior trials. (Verifier correction recorded in the library entry.) & Supports\\
\\[3pt]
Rumble L 2018 \cite{ref4} & Cross-sectional secondary analysis, \seqsplit{multivariate} regression (nationally \seqsplit{representative}, n=6,578) & Women aged 20-24, Indonesia (2012 Indonesian \seqsplit{Demographic} and Health Survey \slash{} Adolescent \seqsplit{Reproductive} Health Survey) & In a nationally \seqsplit{representative} Indonesian sample, \seqsplit{approximately} 17\% of women were married before age 18 and 6\% before age 16; the paper notes research on child marriage in Southeast Asia \seqsplit{specifically} was scarce prior to this nationally \seqsplit{representative} \seqsplit{multivariate} analysis. & Supports\\
\\[3pt]
Williamson HC 2015 \cite{ref10} & RCT (N=130 engaged\slash{}newlywed couples randomized to 1 of 3 \seqsplit{relationship}-education programs, 3-year follow-up) & US engaged and newlywed couples & Couples with acute baseline concerns, including higher physical aggression and alcohol use, benefited LESS from \seqsplit{relationship}-education \seqsplit{intervention} than couples without these concerns; couples with high baseline \seqsplit{satisfaction}\slash{}commitment actually declined faster in \seqsplit{satisfaction} under an intensive skill-based \seqsplit{intervention} versus a low-intensity one. & Supports\\
\\[3pt]
Wilson 2023 \cite{ref99} & instrument-\seqsplit{development}\slash{}validation study & adult \seqsplit{participants} reporting on intimate \seqsplit{relationships} (validation sample) & Reports on validity\slash{}\seqsplit{reliability} of the Demand, Threat, \seqsplit{Surveillance}, and Response-to-Demands subscales of the Coercion in Intimate Partner \seqsplit{Relationships} (CIPR) scale, developed because existing coercive-control measures had documented content-validity gaps (not capturing the full range of demand\slash{}threat\slash{}\seqsplit{surveillance} tactics). & Supports\\
\\[3pt]
Zust BL 2021 \cite{ref118} & Repeated cross-sectional survey (\seqsplit{convenience} sample of Upper-Midwest US \seqsplit{congregations}, 2005 and 2015) & Christian clergy and \seqsplit{congregation} members, US & Faith \seqsplit{commitments} led some domestic-violence victims to endure abuse rather than compromise their faith by leaving; over 10 years, \seqsplit{congregational} support for victims changed slowly; clergy wanted more training in counseling and speaking about domestic violence, and felt \seqsplit{comfortable} making referrals but not \seqsplit{necessarily} providing direct DV counseling themselves. (Verifier correction recorded in the library entry.) & Supports\\
\\[3pt]
Markman HJ 2013 \cite{ref9} & Randomized controlled trial (couples via religious \seqsplit{organizations} randomized to PREP vs. naturally occurring premarital services; n=44\slash{}66\slash{}83) & US couples marrying through religious \seqsplit{organizations}, 8-year follow-up on divorce & No overall difference in divorce rates between PREP (clergy-delivered or \seqsplit{professional}-delivered) and usual premarital services; couples with higher premarital negative-\seqsplit{communication} observed pre-\seqsplit{intervention} were MORE likely to divorce if given PREP, while low-negativity couples benefited; a premarital history of physical aggression showed a similar (marginally \seqsplit{significant}) harmful-moderation pattern. & Complicates\\
\\[3pt]
Miller Elizabeth 2016 \cite{ref110} & Cluster randomized controlled trial (ARCHES) & 4,009 females aged 16-29 across 25 family planning clinics (17 clusters), USA & The universal education\slash{}counseling \seqsplit{intervention} (ARCHES) did NOT \seqsplit{significantly} reduce \seqsplit{reproductive} coercion (ARR 1.50, 95\% CI 0.95-2.35) or partner violence (ARR 1.07, 95\% CI 0.84-1.38) at 1-year follow-up, though it improved secondary outcomes such as \seqsplit{recognition} and self-efficacy. & Complicates\\
\\[3pt]
Robertson 2011 \cite{ref49} & cross-sectional \seqsplit{comparative} study, n=172 (85 men, 87 women), self-report (CTS2 + modified PMWI) & community\slash{}\seqsplit{convenience} samples reporting on own and partner's IPV behavior & Coercive control was associated with IPV similarly for men and women across three samples, and was \seqsplit{predominantly} reciprocal --- both partners commonly reported both receiving and \seqsplit{perpetrating} \seqsplit{controlling} behaviors, regardless of physical-abuse status. & Complicates\\
\\[3pt]
Schneider 2014 \cite{ref47} & \seqsplit{theoretical}\slash{}clinical article (non-empirical) & couples presenting for therapy with \seqsplit{situational} couple violence (SCV) & Argues conjoint treatment can be \seqsplit{appropriate} for \seqsplit{situational} couple violence \seqsplit{specifically} (as distinct from coercive \seqsplit{controlling} violence), proposing an attachment-based safety plan and time-out strategy, while explicitly flagging cautions about \seqsplit{appropriateness} of couples counseling when violence is ongoing. & Complicates\\
\\[3pt]
Portnoy 2020 \cite{srp32791967} & \seqsplit{qualitative} study --- patient interviews (n=10) and provider focus groups (n=29) & US Veterans Health \seqsplit{Administration} patients and providers & Identifies convergent and divergent patient\slash{}provider beliefs about barriers and \seqsplit{facilitators} to screening \seqsplit{specifically} for IPV \seqsplit{perpetration} (not just \seqsplit{victimization}), an area with a much thinner evidence base than \seqsplit{victimization} screening. & Relevant only\\
\\[3pt]
Stern E 2018 \cite{ref114} & \seqsplit{Implementation}\slash{}adaptation review (\seqsplit{qualitative} evaluation and programme \seqsplit{documentation}) & \seqsplit{Indashyikirwa} programme, \seqsplit{implemented} by CARE Rwanda, RWAMREC and RWN, rural Rwanda & \seqsplit{Indashyikirwa} combined a 5-month couples curriculum, sustained community activism by trained couples, women's safe spaces for dedicated support and referral of IPV survivors, and training of opinion leaders to support an enabling \seqsplit{environment}; adaptation was guided by the SASA! fidelity brief. & Relevant only\\
\\[3pt]
\end{longtable}


\subsubsection*{P5 --- Referral quality depends on timeliness, accessibility and follow-through, not offer alone}
% Auto-generated by generate_tex_fragments.py -- do not hand-edit.
\begin{longtable}{@{}L{2.0cm}L{2.0cm}L{2.1cm}L{6.2cm}L{1.9cm}@{}}
\caption{P5: Referral quality depends on timeliness, accessibility and follow-through, not offer alone. Record, design, population/setting, finding as charted (verbatim from the library record), and direction relative to the proposition. n=37 linked records (kg/graph.json edges).}\label{tab:p5}\\
\toprule
\textbf{Record} & \textbf{Design} & \textbf{Population\slash{}setting} & \textbf{Finding (as charted)} & \textbf{Direction}\\
\midrule
\endfirsthead
\multicolumn{5}{l}{\textit{P5 continued}}\\
\toprule
\textbf{Record} & \textbf{Design} & \textbf{Population\slash{}setting} & \textbf{Finding (as charted)} & \textbf{Direction}\\
\midrule
\endhead
\bottomrule
\endfoot
Alghamdi M 2022 \cite{srp34253076} & \seqsplit{Qualitative} dominant design with \seqsplit{supplemental} \seqsplit{quantitative} data (n=8) & Canadian Muslim immigrant women, 6 countries of origin & Pre-migration adverse childhood \seqsplit{experiences} and family history of IPV, combined with post-migration language\slash{}social-isolation barriers and \seqsplit{internalized} \seqsplit{patriarchal} values, shaped help-seeking and the eventual decision to leave abusive \seqsplit{relationships} among 8 Muslim immigrant women. & Supports\\
\\[3pt]
Azalee M 2025 \cite{srp40116416} & \seqsplit{Qualitative} in-depth interviews (n=6) & Health care workers who were IPV survivors, tertiary hospital, Malaysia & Six HCW-survivors used varied coping strategies (seeking help, \seqsplit{spirituality}, avoidance, self-harm, staying\slash{}leaving); survivors relied mainly on coworkers and sought formal help only when situations became critical, hindered by \seqsplit{misconceptions} about IPV, privacy concerns and fear of workplace gossip. & Supports\\
\\[3pt]
Baker CK 2022 \cite{srp34838218} & Systematic review of 34 controlled \seqsplit{quantitative} studies (from 23,902 screened records) & Women \seqsplit{experiencing} IPV, \seqsplit{predominantly} US-based housing\slash{}shelter programmes & [CORRECTION: first\_author is wrong: search lists 'Baker CK \slash{} Sprague S (co-first authorship per abstract group)' but PubMed indexes the first author as Yakubovich (no Baker or Sprague in the author-derived first-author field). Content confirmed: 'screened 23 902 unique records and identified 34 eligible studies'; 'no cumulative evidence of \seqsplit{disadvantages} following any IPV-housing \seqsplit{intervention}. Evidence of benefits was strongest for mental health outcomes, intent to leave partner, perceived safety, and housing and partner-related stress'; 'more research...needed...from outside of the USA.' FIX: change first\_author to Yakubovich.] No cumulative evidence of \seqsplit{disadvantage} from any IPV-housing \seqsplit{intervention}; evidence of benefit was strongest for mental health outcomes, intent to leave partner, perceived safety, and reduced housing\slash{}partner-related stress; most included studies were at high risk of bias, and evidence for long-term\slash{}non-US housing solutions was sparse. & Supports\\
\\[3pt]
Chen 2013 \cite{ref94} & clinical practice review & primary\slash{}family care settings & Describes RADAR (Remember to ask routinely, Ask directly in private, Document, Assess safety, Review options) as a structured clinical tool for \seqsplit{identifying} and responding to IPV, \seqsplit{emphasizing} private \seqsplit{questioning}, safety assessment (weapons, children, \seqsplit{perpetrator} risk), and advocacy-based referral requiring trained links to community resources. & Supports\\
\\[3pt]
Feltner 2018 \cite{srp30357304} & systematic review \seqsplit{commissioned} for USPSTF (global-level guidance body), 30 included studies & adults in healthcare settings, US and comparable settings & [CORRECTION: PMID\slash{}title\slash{}year correct, but the key\_finding materially overstates the abstract's \seqsplit{conclusions} in two ways. (1) The abstract's actual conclusion is: 'trials of IPV screening in adult women did NOT show a reduction in IPV or \seqsplit{improvement} in quality of life over 3 to 18 months' --- screening tools were found to 'reasonably identify' IPV (accuracy claim only, \seqsplit{sensitivity}\slash{}\seqsplit{specificity} varied widely by tool), but the record's claim that the review found 'some \seqsplit{interventions} (e.g., counseling) following screening among women of \seqsplit{reproductive} age reduce IPV' is too broad: only 2 of 11 \seqsplit{intervention} RCTs (1 home-visiting, 1 multi-risk behavioral counseling) showed \seqsplit{significant} reduction, and only in PREGNANT women \seqsplit{specifically}, not the general \seqsplit{reproductive}-age population. (2) 'informing a USPSTF grade-B \seqsplit{recommendation}' conflates this article (the evidence report\slash{}systematic review, PMID 30357304) with the companion USPSTF \seqsplit{Recommendation} Statement, which is a SEPARATE PMID (30357305, also JAMA 2018;320(16):1678-1687) not cited in the record. Fix: rewrite key\_finding to state the evidence report found screening tools moderately accurate but screening ITSELF did not reduce IPV\slash{}improve QoL in trials; only 2 specific \seqsplit{interventions} in pregnant-women subgroups showed benefit; and either cite PMID 30357305 separately for the actual grade-B \seqsplit{recommendation} or explicitly note this record is the evidence report underlying that \seqsplit{recommendation}, not the \seqsplit{recommendation} itself.] Concludes screening \seqsplit{instruments} can accurately identify IPV in health care settings and that some \seqsplit{interventions} (e.g., counseling) following screening among women of \seqsplit{reproductive} age reduce IPV; harms of screening were assessed and found to be low, informing a USPSTF grade-B \seqsplit{recommendation} for \seqsplit{reproductive}-age women. & Supports\\
\\[3pt]
Garcia-Moreno C 2015 \cite{srp25467583} & Series paper \slash{} narrative review with five country case studies & Health systems in low-, middle- and high-income countries & Reviews evidence for clinical \seqsplit{interventions} and health-system components (protocols, training, capacity-building) needed to identify and support women subjected to IPV\slash{}sexual violence following the 2013 WHO guidelines and 67th World Health Assembly resolution; finds \seqsplit{substantial} system and \seqsplit{behavioural} barriers to \seqsplit{implementation}, especially in low- and middle-income countries, and that system \seqsplit{development} has generally been slow. & Supports\\
\\[3pt]
Gerth J 2014 \cite{ref168} & Narrative\slash{}validation review & IPV cases assessed by police, hospital and victim-services staff (multi-\seqsplit{professional}) & [CORRECTION: Actual PubMed title: '[The Ontario Domestic Assault Risk Assessment (ODARA) - Validity und authorised German \seqsplit{translation} of an intimate partner violence screening tool].' Record has 'Validity and authorised German \seqsplit{translation}' --- wrong \seqsplit{conjunction} ('and' vs the German 'und', which PubMed retains in the bilingual title) and drops the subtitle 'of an intimate partner violence screening tool'. FIX: correct the title string. Key finding is accurate: abstract quote 'provide a brief risk assessment tool which can be scored on the basis of only few and easily \seqsplit{collectable} \seqsplit{information}... wider range of \seqsplit{practitioners} (e.g., police officers, hospital and victim services' staff).'] ODARA was designed to be brief and scorable by a wide range of \seqsplit{practitioners} (police, hospital, victim services) using easily \seqsplit{collectable} \seqsplit{information}, to reliably \seqsplit{discriminate} low- from high-risk cases and guide \seqsplit{intervention}. & Supports\\
\\[3pt]
Hill AL 2019 \cite{srp31306331} & Secondary cross-sectional analysis of baseline survey data & 550 sexually active female high school students, USA (school health center attendees, ages 14-19) & [CORRECTION: Prevalence numbers are correct (12\% recent RC, 17\% ARA) among 550 sexually active female students. PROBLEM: key\_finding claims 'both were associated with \seqsplit{differences} in care-seeking and sexual health behaviors' -- but the abstract EXPLICITLY states the opposite for \seqsplit{reproductive} coercion alone: 'There were no observed \seqsplit{significant} \seqsplit{differences} in care-seeking behaviors among those with recent \seqsplit{reproductive} coercion compared with those without.' It is adolescent \seqsplit{relationship} abuse (ARA), not \seqsplit{reproductive} coercion, that was associated with a care-seeking difference (increased odds of STI testing\slash{}treatment, aOR 2.08); the combined ARA+RC-exposed group showed \seqsplit{differences} in sexual-partner\slash{}\seqsplit{contraception} behaviors, not general 'care-seeking'. FIX: rewrite key\_finding to state that RC alone showed NO \seqsplit{significant} care-seeking difference, ARA alone was associated with increased STI-testing care-seeking (aOR 2.08), and the ARA+RC combined group showed elevated odds of specific risk behaviors (older partner, more partners, hormonal-only \seqsplit{contraception}) -- do not state both constructs were 'associated with \seqsplit{differences} in care-seeking'.] 12\% reported recent \seqsplit{reproductive} coercion and 17\% reported adolescent \seqsplit{relationship} abuse; both were associated with \seqsplit{differences} in care-seeking and sexual health behaviors, measured via validated RC and \seqsplit{relationship}-abuse \seqsplit{instruments}. & Supports\\
\\[3pt]
Le V 2025 \cite{srp37649394} & \seqsplit{Qualitative} \seqsplit{implementation} study (16 semi-structured interviews, CFIR-coded) & Catholic diocesan and parish leadership, \seqsplit{Archdiocese} of Chicago Domestic Violence Outreach, USA & [CORRECTION: Key\_finding content (seven CFIR inner-setting constructs and their content) is supported by the abstract. However, first\_author is wrong: record lists 'Le V (\seqsplit{Archdiocese} of Chicago Domestic Violence Outreach study)' but PubMed's actual first author is 'Debinski B'. FIX the first\_author field.] Using the \seqsplit{Consolidated} Framework for \seqsplit{Implementation} Research, seven inner-setting constructs (leadership structure, culture, resources, networking\slash{}\seqsplit{communication}, access to knowledge, \seqsplit{compatibility} with parish values, leadership engagement) were identified as \seqsplit{facilitators} of successful diocesan-level IPV response \seqsplit{implementation}. & Supports\\
\\[3pt]
Lee 2019 \cite{srp31012540} & pre-post quality \seqsplit{improvement} \seqsplit{intervention} study & healthcare providers (nurses, physicians, social workers) at a US clinical site & An EMR-integrated brief IPV screening tool plus provider education plus an automated resource\slash{}referral telephone system was associated with measured \seqsplit{improvement} in provider readiness to screen (assessed via the Domestic Violence Health Care Provider Survey Scale, pre\slash{}post t-test). & Supports\\
\\[3pt]
McCauley HL 2017 \cite{ref17} & Secondary \seqsplit{psychometric} analysis of pooled \seqsplit{multicenter}\slash{}RCT survey data & 4,674 women aged 16-29 in 24 \seqsplit{Pennsylvania} and 5 California family planning clinics, USA & Confirmed a 2-domain structure (pregnancy coercion, condom \seqsplit{manipulation}) and produced a validated 5-item short-form \seqsplit{Reproductive} Coercion Scale (RCS); recent \seqsplit{reproductive} coercion prevalence was 6.3-6.7\% depending on the scale version used. & Supports\\
\\[3pt]
Meiksin R 2020 \cite{srp32460786} & Trial protocol & Women\slash{}girls seeking \seqsplit{contraceptive} services, community-based clinics, Nairobi, Kenya & [CORRECTION: first\_author is wrong: search lists 'Meiksin R' but PubMed indexes the first author as Uysal. Content confirmed: 'ARCHES is a single-session, clinic-based model...Core elements...\seqsplit{opportunity} for patient disclosure of RC and IPV (and subsequent warm referral to local services)' and 'To date, clinic-based GBV \seqsplit{interventions} have not been shown to reduce RC or unintended pregnancy in LMIC settings.' FIX: change first\_author to Uysal.] Describes ARCHES as a single-session, clinic-based enhanced counseling model that creates \seqsplit{opportunity} for patient disclosure of \seqsplit{reproductive} coercion\slash{}IPV followed by warm referral to local services; explicitly notes that to date, no clinic-based GBV \seqsplit{intervention} had been shown to reduce \seqsplit{reproductive} coercion or unintended pregnancy in LMIC settings, motivating this trial. & Supports\\
\\[3pt]
Messing JT 2013 \cite{ref164} & Meta-analytic review of predictive validity (AUC) & Five IPV risk-assessment \seqsplit{instruments} (ODARA, SARA, DA, DVSI, K-SID) across multiple validation studies & ODARA had highest weighted AUC=.666 (k=5), then SARA AUC=.628, DA AUC=.618, DVSI AUC=.582, K-SID AUC=.537; only 9\slash{}20 (45\%) of predictive-validity estimates involved the instrument \seqsplit{administered} correctly per its intended protocol. & Supports\\
\\[3pt]
Nagae M 2004 \cite{ref105} & \seqsplit{Qualitative} interview study & Abused women and health\slash{}legal\slash{}community actors, Khon Kaen province, \seqsplit{northeastern} Thailand & Interviews with 4 abused women and \seqsplit{surrounding} health \seqsplit{professionals}, government\slash{}legal staff and community members in Khon Kaen showed health \seqsplit{professionals} (nurses, community health workers) and, in one case, a community religious leader (a Buddhist monk) played \seqsplit{identifiable} roles in \seqsplit{identifying}, treating and supporting abused women within a multi-actor referral network. & Supports\\
\\[3pt]
O'Doherty L 2015 \cite{ref38} & Cochrane systematic review and meta-analysis of 13 RCTs\slash{}quasi-RCTs (n=14,959) & Women in antenatal, women's health, emergency department and primary care settings, mostly high-income urban settings & Screening increased clinical \seqsplit{identification} of IPV victims (OR 2.95, 95\% CI 1.79-4.87, n=10,074) but there was no evidence of effect on referrals to support services (OR 2.24, 95\% CI 0.64-7.86, only 2 studies, n=1298), no evidence screening reduced women's re-exposure to violence at 3-18 months, and \seqsplit{insufficient} evidence on health outcomes or harm beyond the immediate post-screening period; authors conclude screening alone is \seqsplit{insufficiently} supported to justify universal \seqsplit{implementation} absent stronger referral\slash{}outcome evidence. & Supports\\
\\[3pt]
Sharifnia A 2025 \cite{srp39425490} & Meta-\seqsplit{ethnography} (33 \seqsplit{qualitative} studies \seqsplit{synthesized}), 6 databases searched & Muslim women, global (multi-country \seqsplit{qualitative} evidence) & Synthesis of 33 \seqsplit{qualitative} studies identified four major themes - experience\slash{}impact of abuse, risk factors, help-seeking, and coping - and found religion functions both as a risk factor for abuse and a barrier to help-seeking, and \seqsplit{simultaneously} as a coping resource for some women. & Supports\\
\\[3pt]
Sprague 2017 \cite{srp28649297} & scoping review, 43 included studies & health care settings, multiple countries & Of 43 studies across 9 categories of IPV assistance programmes (\seqsplit{counselling}\slash{}advocacy, safety planning, referral, resource provision, home visitation, case management, videos, provider cueing, system changes), programmes using a single active-assistance type combined with a counsellor\slash{}advocate reported positive results in 77.8\% of studies, suggesting simpler, advocate-led \seqsplit{interventions} outperform complex bundled ones. & Supports\\
\\[3pt]
Trabold N 2014 \cite{ref109} & 3-year community-based \seqsplit{participatory} mixed-methods \seqsplit{demonstration} project across 4 community health centers & Urban, \seqsplit{predominantly} indigent safety-net primary care population, USA & [CORRECTION: first\_author is wrong: search lists 'Trabold N' but PubMed indexes the first author as Rhodes. Design \seqsplit{description} matches ('3-year community-based mixed-method \seqsplit{participatory} research project involving four community health centers...\seqsplit{predominately} indigent urban population'). key\_finding's general claim of '\seqsplit{significant} system-level challenges' is supported at a high level ('Results highlight the challenges, but also the \seqsplit{opportunities}, for \seqsplit{implementing} the new USPSTF guidelines...in resource-limited primary care settings') but the abstract itself does not itemize 'workflow, staff capacity, referral \seqsplit{infrastructure}' as specific named barriers -- that level of detail is not verifiable from the abstract alone (may be in full text). FIX: change first\_author to Rhodes; consider softening the \seqsplit{parenthetical} specifics or citing they come from the full text, not the abstract.] \seqsplit{Implementing} USPSTF-guideline-consistent universal IPV screening and referral in resource-limited urban primary care settings faced \seqsplit{significant} system-level challenges (workflow, staff capacity, referral \seqsplit{infrastructure}), even with a dedicated multi-year system-level \seqsplit{intervention} designed to raise detection and response. & Supports\\
\\[3pt]
US Preventive Services Task Force 2025 \cite{ref40} & guideline \slash{} \seqsplit{recommendation} statement (systematic-review-based) & US \seqsplit{adolescents} and adults, pregnant\slash{}postpartum and \seqsplit{reproductive}-age women; older\slash{}vulnerable adults without recognized signs of abuse & USPSTF concludes screening for IPV in women of \seqsplit{reproductive} age (including pregnant\slash{}postpartum), combined with provision or referral to \seqsplit{multicomponent} \seqsplit{interventions}, has moderate net benefit (B \seqsplit{recommendation}). For caregiver abuse\slash{}neglect of older or vulnerable adults, benefits and harms remain \seqsplit{undetermined} (I statement) even in this 2025 update. & Supports\\
\\[3pt]
Waalen J(CFPSA 2014 \cite{ref111} & Systematic review of 9 RCTs of physician-directed domestic violence education & \seqsplit{Postgraduate} trainee and practising physicians & [CORRECTION: first\_author is wrong: search lists 'Waalen J' but PubMed indexes the first author as Zaher. Content confirmed precisely: 'Brief \seqsplit{interventions} for \seqsplit{postgraduate} trainee physicians improved knowledge but did not seem to affect behaviour...Physician training combined with system support \seqsplit{interventions} seemed to benefit domestic violence victims and increase referrals to domestic violence support resources.' FIX: change first\_author to Zaher.] Brief \seqsplit{educational} \seqsplit{interventions} for \seqsplit{postgraduate} trainees improved knowledge but did NOT change screening\slash{}\seqsplit{identification} behaviour; only \seqsplit{multifaceted} training combined with system-support \seqsplit{interventions} (reminders, access to victim-support referral pathways, audit\slash{}feedback) was associated with increased referrals to domestic-violence support resources. & Supports\\
\\[3pt]
Wang PL 2015 \cite{ref165} & Validation cohort study & 543 female IPV victims in a Taiwanese victim-services program & TIPVDA showed strong \seqsplit{discriminant} power and ROC-AUC support for predicting lethal-danger risk, with stronger predictive power at the high-\seqsplit{dangerousness} end. & Supports\\
\\[3pt]
Williamson HC 2014 \cite{ref106} & Cross-sectional survey study (N=2,126 married \seqsplit{individuals}) & US married adults, general population sample & Contrary to the assumption that premarital education reduces later need for counseling, receiving premarital education covaried with an INCREASED likelihood of later couples counseling (a 'gateway' effect), with the \seqsplit{association} stronger among African Americans and those with lower income\slash{}education. & Supports\\
\\[3pt]
Williamson HC 2019 \cite{ref107} & Cross-sectional self-report study (N=231 ethnically diverse newlywed couples in low-income \seqsplit{neighborhoods}) & US low-income newlywed couples & Top perceived barriers to seeking couple therapy were cost and \seqsplit{uncertainty} about where to go for help; direct experience with premarital education (or indirect experience via a social-network member's couple therapy) was associated with reduced barriers and greater likelihood of \seqsplit{subsequently} seeking therapy. & Supports\\
\\[3pt]
Zust BL 2021 \cite{ref118} & Repeated cross-sectional survey (\seqsplit{convenience} sample of Upper-Midwest US \seqsplit{congregations}, 2005 and 2015) & Christian clergy and \seqsplit{congregation} members, US & [CORRECTION: Key\_finding is well supported: 'Victims find strength in their faith and would rather endure the violence at all costs to keep a family or a marriage together, than to compromise their faith by leaving' and 'Clergy felt \seqsplit{comfortable} in making referrals for \seqsplit{professional} counseling.' However relevance\_to\_tphe overstates the referral point: the abstract's very next clause is 'while the majority of members would prefer counseling with their pastor if they were in a violent \seqsplit{relationship}' --- i.e. \seqsplit{congregants} actually want direct pastoral counseling, not referral, and clergy said they want MORE training to counsel\slash{}speak about DV themselves, not that they prefer referral over direct counseling. FIX: change 'clergy themselves report needing referral pathways rather than attempting DV counseling directly' to something like 'clergy report comfort in making referrals, while members\slash{}clergy \seqsplit{simultaneously} want more direct pastoral counseling capacity --- a live tension rather than clean support for a refer-not-counsel policy that T-PHE's safeguard should note, not paper over.'] Faith \seqsplit{commitments} led some domestic-violence victims to endure abuse rather than compromise their faith by leaving; over 10 years, \seqsplit{congregational} support for victims changed slowly; clergy wanted more training in counseling and speaking about domestic violence, and felt \seqsplit{comfortable} making referrals but not \seqsplit{necessarily} providing direct DV counseling themselves. & Supports\\
\\[3pt]
\seqsplit{Grisurapong} S 2002 \cite{ref104} & \seqsplit{Intervention}\slash{}comparison case-site \seqsplit{description} & Women \seqsplit{experiencing} physical or sexual assault, Khon Kaen provincial hospital, Thailand & Describes the \seqsplit{establishment} and early operation of a hospital-based One Stop Crisis Center (OSCC) for women survivors of violence in Khon Kaen, Thailand, compared against a non-\seqsplit{intervention} provincial hospital site. & Supports\\
\\[3pt]
Hogan NR 2024 \cite{ref171} & Cross-sectional analysis, n=190 & \seqsplit{Individuals} with IPV histories assessed via SARA-V3, Canada (Indigenous vs non-Indigenous) & Structured-\seqsplit{professional}-judgment risk ratings (not only actuarial tools) showed racial \seqsplit{disparities} between Indigenous and non-Indigenous \seqsplit{participants} in risk factors and summary ratings. & Complicates\\
\\[3pt]
Lee Y 2019 \cite{srp29384426} & Randomized controlled pilot trial (online modules, baseline and 3-month follow-up) & Korean American faith leaders, \seqsplit{southeastern} USA (n=55; \seqsplit{intervention} n=27, control n=28) & The \seqsplit{intervention} group showed \seqsplit{significant} \seqsplit{improvement} in knowledge of IPV resources and attitudes against IPV compared to control, but gains in self-efficacy and prevention\slash{}\seqsplit{intervention} behaviours, while greater than control, were not \seqsplit{statistically} \seqsplit{significant}. & Complicates\\
\\[3pt]
Onukwugha FI 2025 \cite{ref108} & Mixed-methods \seqsplit{implementation} study embedded in a matched-pair cluster-controlled trial (75 providers, 20 facilities) & Antenatal care and family planning clinics, Nigeria; nurse-midwives and community health extension workers trained in WHO LIVES + ARCHES & [CORRECTION: first\_author is wrong: search lists 'Onukwugha FI' but PubMed indexes the first author as Oduenyi. Content is otherwise well supported: 'Quality assurance standards increased over time (P $\le$ 0.05)'; 'average scores of \seqsplit{approximately} 3 out of a maximum of 4'; 'providers initially perceived the \seqsplit{intervention} as additional work...lack of funds and limited staff resources were found to be key barriers'; 'pulse survey showed no \seqsplit{significant} change over time for \seqsplit{feasibility}, \seqsplit{appropriateness} and \seqsplit{acceptability}.' FIX: change first\_author to Oduenyi.] Quality-assurance readiness scores for GBV first-line response increased \seqsplit{significantly} over 9 months (P<=0.05); providers rated \seqsplit{integration} as \seqsplit{appropriate}\slash{}acceptable (\textasciitilde{}3 of 4 on average) but reported it initially felt like added work, and lack of funds\slash{}staff was a key barrier; \seqsplit{feasibility}\slash{}\seqsplit{acceptability} pulse scores showed no \seqsplit{significant} change over time despite improved facility readiness. & Complicates\\
\\[3pt]
US Preventive Services Task Force 2018 \cite{srp30357305} & guideline \slash{} \seqsplit{recommendation} statement & US women of \seqsplit{reproductive} age; older\slash{}vulnerable adults & With moderate certainty, screening plus referral to ongoing support services has moderate net benefit; explicitly states evidence does NOT support the \seqsplit{effectiveness} of brief \seqsplit{interventions} or merely providing referral \seqsplit{information} in the absence of ongoing supportive \seqsplit{intervention} components, and that benefit evidence is \seqsplit{concentrated} in pregnant\slash{}postpartum \seqsplit{populations}. & Complicates\\
\\[3pt]
Bundock Kerrie 2020 \cite{srp29683049} & Systematic review (19 studies) & Adolescent victims of dating violence, mixed \seqsplit{international} samples (English-language literature) & \seqsplit{Adolescents} \seqsplit{overwhelmingly} favored informal help sources (friends most common) over formal services; females were generally more likely than males to seek help (with \seqsplit{inconsistencies}); \seqsplit{measurement} and \seqsplit{definitions} of both violence and help-seeking varied \seqsplit{substantially} across included studies, limiting \seqsplit{comparability}. & Relevant only\\
\\[3pt]
Gutowski 2023 \cite{srp35611345} & instrument-\seqsplit{development} study, n=222 survivor-mothers, \seqsplit{exploratory} factor analysis + Rasch analysis & US survivor-mothers involved in family law \seqsplit{proceedings} & Developed a 14-item, two-subscale Legal Abuse Scale (LAS) to measure how family-court and legal processes can be misused by partners to continue coercive control post-separation. & Relevant only\\
\\[3pt]
Hameed M 2020 \cite{srp32608505} & Cochrane systematic review and meta-analysis of 33 RCTs (n=5517 women) & Women aged 16+ with recent or lifetime IPV experience, mostly high-income countries, recruited from healthcare, community, shelter\slash{}refuge settings & \seqsplit{Psychological} therapies probably reduce depression at medium-term follow-up (SMD -0.24, 95\% CI -0.47 to -0.01; 4 trials, 600 women, moderate certainty) and may reduce anxiety, but evidence is uncertain or absent for self-efficacy, PTSD, re-exposure to IPV, and safety planning; harm data are limited (only 1 of 33 trials used a validated harm-outcome scale). & Relevant only\\
\\[3pt]
Michaels-Igbokwe C 2016 \cite{srp26924488} & Economic evaluation alongside a cluster randomized controlled trial & SASA! \seqsplit{intervention} \seqsplit{communities}, Kampala, Uganda, provider \seqsplit{perspective}, 4 years of \seqsplit{programming} & Total \seqsplit{intervention} cost over 4 years was estimated at US\$553,252, averaging US\$5 per person reached annually; an estimated 1,201 cases of past-year physical IPV were averted (90\% CI 97-2,307) at a cost of US\$460 per case averted. & Relevant only\\
\\[3pt]
Schenck D 2018 \cite{ref160} & Ethical analysis \slash{} case-based policy proposal & Domino liver donor transplant candidates (pediatric case example) & [CORRECTION: Abstract quote: 'We propose a Domino Donor Advocate-based on the concept of the \seqsplit{independent} living donor advocate... We discuss the need for 2 distinct consent processes separated in time to ensure that potential domino donors... give a truly voluntary consent.' Key finding accurate. Transplant-medicine analogy again phrased as 'supports R4 no-stake assessor and P5 referral timing' --- FIX to analogy language per \seqsplit{instructions}.] Proposes a 'Domino Donor Advocate', modeled explicitly on the \seqsplit{independent} living donor advocate, plus two consent processes separated in time, to protect \seqsplit{voluntariness} when the same party stands to gain from securing consent. & Relevant only\\
\\[3pt]
Stern E 2018 \cite{ref114} & \seqsplit{Implementation}\slash{}adaptation review (\seqsplit{qualitative} evaluation and programme \seqsplit{documentation}) & \seqsplit{Indashyikirwa} programme, \seqsplit{implemented} by CARE Rwanda, RWAMREC and RWN, rural Rwanda & \seqsplit{Indashyikirwa} combined a 5-month couples curriculum, sustained community activism by trained couples, women's safe spaces for dedicated support and referral of IPV survivors, and training of opinion leaders to support an enabling \seqsplit{environment}; adaptation was guided by the SASA! fidelity brief. & Relevant only\\
\\[3pt]
World Health \seqsplit{Organization} 2013 (no PMID) & Global clinical and policy guideline & Global health-system guidance for responding to women subjected to IPV or sexual violence & \seqsplit{Establishes} WHO's first-line support principles for clinicians (the basis of the later-named LIVES model: Listen, Inquire about needs, Validate, Enhance safety, Support\slash{}refer), recommends against mandatory universal screening in the absence of adequate response systems, and sets minimum \seqsplit{requirements} (privacy, \seqsplit{confidentiality}, trained staff, referral pathways) before any clinical inquiry about violence should occur. & Relevant only\\
\\[3pt]
World Health \seqsplit{Organization} 2014 (no PMID) & Global clinical practice handbook \seqsplit{operationalizing} the 2013 WHO guidelines & Frontline health workers globally, first-line clinical response to women disclosing IPV or sexual violence & \seqsplit{Operationalizes} the LIVES first-line support protocol (Listen, Inquire, Validate, Enhance safety, Support) into a step-by-step clinical tool for frontline providers, with explicit guidance on \seqsplit{confidentiality}, \seqsplit{documentation}, safety planning, and structured referral rather than passive \seqsplit{information} provision. & Relevant only\\
\\[3pt]
\end{longtable}


\subsubsection*{P6 --- Conventional programme success can coexist with \tphe{} failure}
% Auto-generated by generate_tex_fragments.py -- do not hand-edit.
\begin{longtable}{@{}L{2.0cm}L{2.0cm}L{2.1cm}L{6.2cm}L{1.9cm}@{}}
\caption{P6: Conventional programme success can coexist with T-PHE failure. Record, design, population/setting, finding as charted (verbatim from the library record), and direction relative to the proposition. n=6 linked records (kg/graph.json edges).}\label{tab:p6}\\
\toprule
\textbf{Record} & \textbf{Design} & \textbf{Population\slash{}setting} & \textbf{Finding (as charted)} & \textbf{Direction}\\
\midrule
\endfirsthead
\multicolumn{5}{l}{\textit{P6 continued}}\\
\toprule
\textbf{Record} & \textbf{Design} & \textbf{Population\slash{}setting} & \textbf{Finding (as charted)} & \textbf{Direction}\\
\midrule
\endhead
\bottomrule
\endfoot
Fauk NK 2021 \cite{ref103} & \seqsplit{Qualitative} in-depth interviews (n=92: 52 women, 40 men) & People living with HIV, Belu and Yogyakarta, Indonesia (Islamic and Javanese cultural context) & Cultural practices (bride wealth, spousal-dispute management) and Islamic religious beliefs about husband-wife \seqsplit{relationships}, forbidden premarital sex, and condom use were reported by \seqsplit{participants} as shaping spousal sexual behaviours that \seqsplit{contributed} to HIV \seqsplit{transmission} within marriage. & Supports\\
\\[3pt]
Islam MK 2021 \cite{ref5} & Cross-sectional survey, \seqsplit{multivariable} logistic regression (n=486) & Rohingya refugee women, Cox's Bazar settlement, Bangladesh & 61.32\% of surveyed women had \seqsplit{experienced} child marriage; child marriage was associated with \seqsplit{significantly} higher odds of justifying husband-\seqsplit{perpetrated} physical abuse (AOR=2.71, 95\%CI 1.78-4.11) and of \seqsplit{experiencing} IPV in the past 12 months (AOR=1.72, 95\%CI 1.13-2.61). & Supports\\
\\[3pt]
Jacubowski N 2008 \cite{ref112} & \seqsplit{Documentary}\slash{}archival analysis and expert-informant interviews & Married women, Indonesia (Muslim-majority context) & \seqsplit{Traditional} practices including polygamy, early marriage and contract marriage (mut'a) were identified as enhancing married Indonesian women's likelihood of HIV \seqsplit{acquisition} within marriage, \seqsplit{challenging} the assumption that marital status is protective. & Supports\\
\\[3pt]
Hawkins AJ 2008 \cite{ref6} & Meta-analysis (86 reports, 117 studies, >500 effect sizes) & Mostly US\slash{}Western MRE program \seqsplit{participants}; \seqsplit{predominantly} White, middle-class samples & \seqsplit{Relationship}-quality effect sizes for \seqsplit{experimental} studies ranged d=.30-.36 and \seqsplit{communication}-skills effects d=.43-.45; moderate-dosage programs \seqsplit{outperformed} low-dosage; the authors flag that racial\slash{}ethnic and economic diversity was too thin to draw reliable \seqsplit{conclusions} for \seqsplit{disadvantaged} couples, and that outcomes \seqsplit{policymakers} care about (\seqsplit{relationship} stability, aggression) were rarely measured. & Complicates\\
\\[3pt]
Phonkacha N 2026 \cite{srp42300599} & Cross-sectional national survey, binary probit regression & General population, national survey, Thailand (Muslim identity modeled as a covariate) & Muslim identity, poverty, and residence in southern\slash{}central Thailand correlated with higher IPV tolerance for both genders; internet access (as distinct from device ownership) \seqsplit{significantly} reduced IPV tolerance, suggesting a role for \seqsplit{information} access in shifting attitudes. & Complicates\\
\\[3pt]
Sayers SL 1998 \cite{srp9779330} & Narrative review of 13 program-evaluation studies & General review of \seqsplit{behaviorally} oriented marriage-\seqsplit{preparation} programs, mostly US & \seqsplit{Behaviorally} oriented, skills-based marriage-\seqsplit{preparation} programs across 13 reviewed studies showed evidence of producing behavioral changes that may help prevent marital \seqsplit{dysfunction}, though program-evaluation \seqsplit{methodology} at the time was described as only modestly improved and needing further rigor. & Relevant only\\
\\[3pt]
\end{longtable}


\subsubsection*{P7 --- Repairability is necessary because relevant information emerges over time}
% Auto-generated by generate_tex_fragments.py -- do not hand-edit.
\begin{longtable}{@{}L{2.0cm}L{2.0cm}L{2.1cm}L{6.2cm}L{1.9cm}@{}}
\caption{P7: Repairability is necessary because relevant information emerges over time. Record, design, population/setting, finding as charted (verbatim from the library record), and direction relative to the proposition. n=6 linked records (kg/graph.json edges).}\label{tab:p7}\\
\toprule
\textbf{Record} & \textbf{Design} & \textbf{Population\slash{}setting} & \textbf{Finding (as charted)} & \textbf{Direction}\\
\midrule
\endfirsthead
\multicolumn{5}{l}{\textit{P7 continued}}\\
\toprule
\textbf{Record} & \textbf{Design} & \textbf{Population\slash{}setting} & \textbf{Finding (as charted)} & \textbf{Direction}\\
\midrule
\endhead
\bottomrule
\endfoot
Michau L 2021 \cite{ref113} & Programme-theory\slash{}programme-\seqsplit{description} paper (non-empirical evaluation article) & Raising Voices programme learning, Uganda-originated model adapted for broader use & The paper describes the rationale for revising SASA! after ten years of programme learning and research, detailing how the revised 'SASA! Together' model \seqsplit{incorporated} field learning, and reflects on debates such as how best to scale up violence-prevention programmes. & Supports\\
\\[3pt]
Silliman B 1999 \cite{ref116} & Narrative review & General review of marriage-\seqsplit{preparation} program practice, US & Reviews two decades of research on marriage-\seqsplit{preparation} programs and gives design \seqsplit{recommendations} covering needs assessment, timing\slash{}duration, format, content, provider \seqsplit{characteristics}, and promotion; does not report new outcome data. & Supports\\
\\[3pt]
Subramanee SD 2022 \cite{srp36429857} & Systematic review (PRISMA, 13 included studies of 520 screened), PROSPERO \seqsplit{CRD42020190410} & South Asian \seqsplit{populations} (India, Bangladesh, Pakistan, Nepal, Sri Lanka), including Muslim subgroups & Consistent correlates of child marriage were rural residence, low education, poor economic background, low mass-media exposure and religion (Hindu and Muslim in particular countries); downstream maternal-health effects included lower antenatal care use, \seqsplit{institutional} delivery and skilled birth attendance. & Supports\\
\\[3pt]
Upadhyay UD 2014 \cite{ref16} & Instrument \seqsplit{development} and validation (cross-sectional survey, factor analysis) & 1,892 women at 13 family planning and 6 abortion facilities, USA & Generated 26 candidate items, retained 14 items across 3 subscales (freedom from coercion, \seqsplit{communication}, decision-making) via factor analysis; the freedom-from-coercion and \seqsplit{communication} subscales were inversely associated with \seqsplit{unprotected} sex in the past 3 months, supporting construct validity. & Supports\\
\\[3pt]
Clyde TL 2020 \cite{ref115} & Narrative review\slash{}position paper & General discussion of \seqsplit{contemporary} premarital-education practice, US-centric & Reviews social trends (\seqsplit{individualism}\slash{}commitment \seqsplit{ambivalence}, changing marriage attitudes, premarital \seqsplit{relationship} histories, media \seqsplit{environment}) and argues current premarital \seqsplit{interventions} need updated content and structure to remain relevant to engaged couples today; offers four proposals for \seqsplit{redesigning} premarital \seqsplit{interventions}. & Relevant only\\
\\[3pt]
Upadhyay UD 2021 \cite{ref18} & Instrument \seqsplit{development} and validation (\seqsplit{qualitative} item generation + national survey, EFA, logistic regression) & 1,117 sexually active young people aged 15-24, USA (national sample) & Developed a 23-item, 7-subscale Sexual and \seqsplit{Reproductive} \seqsplit{Empowerment} Scale for \seqsplit{Adolescents}\slash{}Young Adults (including 'choice of partners, marriage, and children' as a named subscale); validation via logistic regression supported construct validity. & Relevant only\\
\\[3pt]
\end{longtable}


\subsection{Evidence gaps}
Table~\ref{tab:gaps} excerpts the \texttt{gaps\_noted} statements logged by each search pass at the
time it was run, grouped into the four themes named in this review's scope: Thai/Thai-Muslim
samples, premarital-specific trials, decision-control instruments, and Layer~0/Layer~2 evidence
(\tphe{}'s two proposed, not-yet-accepted extensions). These are the search log's own statements
about what it did not find, condensed to fit the table column where needed, not a new claim
authored for this review.

% Auto-generated by generate_tex_fragments.py -- do not hand-edit.
\begin{longtable}{@{}L{3.0cm}L{2.4cm}L{9.0cm}@{}}
\caption{Evidence gaps, excerpted from the gaps\_noted field logged in each search-*.json angle file at the time of the search pass. Text is the search log's own statement, condensed to fit the table column, not a new claim.}\label{tab:gaps}\\
\toprule
\textbf{Gap theme} & \textbf{Search angle} & \textbf{Logged gap statement}\\
\midrule
\endfirsthead
\multicolumn{3}{l}{\textit{Table \thetable{} continued}}\\
\toprule
\textbf{Gap theme} & \textbf{Search angle} & \textbf{Logged gap statement}\\
\midrule
\endhead
\bottomrule
\endfoot
Thai \slash{} Thai-Muslim samples & assessor-\seqsplit{independence}-checkin-models & No Thailand-specific or premarital-education-specific record was found for this angle; angle (a)-(d) evidence is drawn from adjacent fields (organ \seqsplit{transplantation} ethics, IPV risk assessment, suicide\slash{}self-harm brief-contact trials, US pediatric\slash{}prevention-science \seqsplit{implementation} studies) and must be explicitly flagged as analogical transfer, not direct T-PHE evidence, in any manuscript use.\\
\\[3pt]
Thai \slash{} Thai-Muslim samples & couple-work-safety-coercive-control & No PubMed-indexed study was found evaluating a couple-based premarital or pre-nikah education programme with an embedded coercive-control screening\slash{}stop-rule mechanism in a Muslim or Thai\slash{}Southeast Asian \seqsplit{intercultural} setting --- the exact population T-PHE targets.\\
\\[3pt]
Thai \slash{} Thai-Muslim samples & decision-autonomy-measures & No PubMed-indexed decision-autonomy or coercion-\seqsplit{recognition} instrument was found validated \seqsplit{specifically} in a Thai population, Thai Muslim population, or Thai \seqsplit{intercultural}-couple population; the closest cultural\slash{}religious matches found (Türkiye, Lebanon, Rohingya\slash{}Syrian refugees) are \seqsplit{geographically} and culturally distant from Thailand.\\
\\[3pt]
Thai \slash{} Thai-Muslim samples & decision-autonomy-measures & Cross-cultural adaptation studies located here show mixed results (Hazar 2024 Turkey: acceptable but uneven subscale \seqsplit{reliability}; Dagher 2024 Lebanon: sub-optimal CFA fit even after item removal) --- this is direct evidence against assuming any existing decision-control\slash{}autonomy instrument would transfer cleanly to a new Thai Muslim\slash{}\seqsplit{intercultural} context without \seqsplit{independent} local validation, a gap the T-PHE abstract...\\
\\[3pt]
Thai \slash{} Thai-Muslim samples & digital-self-assessment-ipv & No Thailand-specific RCT or \seqsplit{effectiveness} (as opposed to \seqsplit{preliminary} expert-evaluation) data were found for a scored self-check tool; Chuemchit 2026 (PMID 42196762) is \seqsplit{development}\slash{}heuristic-evaluation stage only, not yet tested with the target population or measured for \seqsplit{behavioural}\slash{}health outcomes.\\
\\[3pt]
Thai \slash{} Thai-Muslim samples & humility-compassion-self-regulation & No Muslim-sample or Thailand-based record was found for the \seqsplit{intellectual}-humility or self-compassion \seqsplit{literatures} (angles a\slash{}c); the \seqsplit{religiosity}-IPV evidence located (36331229, 33704630) is Brazilian-mixed-sample and Christian-family literature \seqsplit{respectively}, not Muslim samples, despite the query explicitly requesting Muslim samples where indexed --- none were indexed for this specific construct pairing.\\
\\[3pt]
Thai \slash{} Thai-Muslim samples & ipv-screening-response & No PubMed-indexed evidence located on IPV screening\slash{}response \seqsplit{specifically} among Muslim or Thai\slash{}Southeast Asian couples; the Nigeria trials (Sokoto, Ebonyi) are the closest Muslim-population-adjacent evidence but are clinical, not premarital-education, settings and not Southeast Asian.\\
\\[3pt]
Thai \slash{} Thai-Muslim samples & islamic-medical-ethics-tier-one & No PubMed-indexed record \seqsplit{specifically} evaluates premarital health education content (not screening, not marriage-age demography) delivered through an explicitly named Islamic ethical frame (faqr\slash{}taqwa\slash{}amanah\slash{}'adl\slash{}shura) in a Thai or Southeast Asian Muslim population --- the closest analogues found (Bangladesh mosque-based diabetes RCT, PMID 41114978; Somali refugee mosque-based trauma RCT, PMID 39186273) are faith-i...\\
\\[3pt]
Thai \slash{} Thai-Muslim samples & islamic-medical-ethics-tier-one & No BMC Medical Ethics systematic review specific to Islamic bioethics was found; the one Islamic-adjacent BMC Medical Ethics record located (PMID 41491519, a scoping review on social egg freezing) was excluded as \seqsplit{insufficiently} on-topic (not focused on a Muslim population or the five named ethical terms).\\
\\[3pt]
Thai \slash{} Thai-Muslim samples & \seqsplit{manipulation}-literacy-\seqsplit{inoculation} & No PubMed-indexed Thailand-resident (as opposed to diaspora) study on \seqsplit{recognition} of emotional\slash{}\seqsplit{psychological} abuse or coercive control was found; the two most population-relevant records (26289454, 26442989) are Canada\slash{}US-diaspora South\slash{}Southeast Asian Muslim and Southeast Asian samples, not Thailand-resident \seqsplit{populations} -- a genuine evidence gap for angle (e) as specified.\\
\\[3pt]
Thai \slash{} Thai-Muslim samples & muslim-thailand-southeast-asia & No PubMed-indexed evaluation study of an Islamic\slash{}faith-based premarital education or marriage-\seqsplit{preparation} course (content, uptake, or outcomes) specific to Thailand, Malaysia, or Indonesia was found; premarital-course literature retrieved was \seqsplit{overwhelmingly} about genetic\slash{}\seqsplit{thalassaemia} premarital screening, not relational or \seqsplit{safeguarding} education.\\
\\[3pt]
Thai \slash{} Thai-Muslim samples & muslim-thailand-southeast-asia & No PubMed-indexed study on secret, \seqsplit{unregistered}, or informal (nikah siri \slash{} nikah tidak dicatatkan-type) religious marriage prevalence, drivers, or \seqsplit{consequences} among Muslim \seqsplit{populations} in Thailand\slash{}Malaysia\slash{}Indonesia was located via these queries.\\
\\[3pt]
Thai \slash{} Thai-Muslim samples & muslim-thailand-southeast-asia & No PubMed-indexed study on Thailand's Deep South (Pattani\slash{}Yala\slash{}Narathiwat) Muslim population \seqsplit{specifically} addressing IPV, forced marriage, or premarital programmes was found; the Thailand-specific IPV\slash{}OSCC literature retrieved (\seqsplit{Grisurapong} 2002, Nagae 2004, Phonkacha 2026) is either general-population or historical, not Muslim-population-specific within Thailand.\\
\\[3pt]
Thai \slash{} Thai-Muslim samples & muslim-thailand-southeast-asia & Evidence on Muslim \seqsplit{populations} in Malaysia and Indonesia specific to premarital decision-control, forced-marriage refusal, or STOP-RULE-type joint-versus-separated counseling protocols was not found in PubMed; most Southeast Asian Muslim marriage literature indexed in PubMed concerns child-marriage prevalence\slash{}\seqsplit{determinants} (largely Indonesia) or HIV \seqsplit{vulnerability} within marriage, not programme design or \seqsplit{safeguardin}...\\
\\[3pt]
Thai \slash{} Thai-Muslim samples & muslim-thailand-southeast-asia & Rohingya evidence, while ethnically\slash{}\seqsplit{religiously} proximate and \seqsplit{geographically} linked to the Bay of Bengal\slash{}Southeast Asia border region, is drawn from a Bangladesh refugee-camp setting, not Thailand proper; it is included as the closest available regional Muslim-population IPV\slash{}child-marriage evidence but should not be read as directly \seqsplit{representative} of Thai Muslim or \seqsplit{intercultural} couples.\\
\\[3pt]
Thai \slash{} Thai-Muslim samples & peer-online-support-digital-safe-space & No PubMed-indexed record was found specific to Thailand, or to a Thai Muslim population, on digital peer support, online IPV support, or moderated online \seqsplit{communities}; the Muslim-specific record found (PMID 35536616) is set in Australia, and other Muslim\slash{}health hits from the search (PMID 37928060 anti-Muslim \seqsplit{discrimination}, PMID 38040423 Sikh\slash{}Muslim kidney transplant views, PMID 31120322 US Muslim \seqsplit{contraceptive} bel...\\
\\[3pt]
Thai \slash{} Thai-Muslim samples & premarital-education-outcomes & No Thailand-specific or Southeast Asian premarital\slash{}marriage education outcome study was located; the closest cultural analogues found are Iranian (PMID 40783649), Turkish (PMID 42442188), Saudi (PMID 36620784) and South Korean (PMID 39757299) studies -- none are Southeast Asian or \seqsplit{specifically} Thai Muslim\slash{}\seqsplit{intercultural} \seqsplit{populations}.\\
\\[3pt]
Thai \slash{} Thai-Muslim samples & primary-prevention-frameworks & No papers were found \seqsplit{specifically} evaluating IPV-prevention programme mechanisms among Muslim or \seqsplit{intercultural} couples in Thailand or Southeast Asia; the Nigerian Muslim\slash{}Christian \seqsplit{congregational} trial (PMID 37329160) is the nearest cross-cultural, mixed-faith analogue found.\\
\\[3pt]
Thai \slash{} Thai-Muslim samples & primary-prevention-frameworks & Faith-based \seqsplit{intervention} evidence located skews toward Christian\slash{}Catholic and one Korean-immigrant Christian clergy sample plus one Nigerian Muslim\slash{}Christian trial; \seqsplit{specifically} Islamic religious-authority-led premarital education \seqsplit{evaluations} were not located in this search.\\
\\[3pt]
Premarital-specific trials & assessor-\seqsplit{independence}-checkin-models & No PubMed-indexed literature was found that directly studies 'conflict of interest of a wedding-service provider acting as premarital-education assessor' --- this exact \seqsplit{configuration} does not appear to be a studied topic; evidence is \seqsplit{necessarily} analogical (organ-donor advocacy, clinical dual-\seqsplit{relationship} ethics, facial-plastic-surgery COI) rather than domain-matched.\\
\\[3pt]
Premarital-specific trials & couple-work-safety-coercive-control & Coercive-control \seqsplit{measurement} \seqsplit{instruments} identified (CIPR, C-CAS, LAS, C-CBS-R, CASR-SF) are validated \seqsplit{predominantly} in Western (US, Australian, Canadian) or Chinese-community samples; only one (Tiwari 2015, Chinese sample) offers a non-Western cross-cultural precedent, and none has been validated for premarital screening \seqsplit{specifically} (all are post-hoc\slash{}\seqsplit{retrospective} abuse measures, not \seqsplit{prospective} risk-of-future-c...\\
\\[3pt]
Premarital-specific trials & couple-work-safety-coercive-control & No cost-\seqsplit{effectiveness}, \seqsplit{implementation}-fidelity, or long-term (multi-year) outcome data were found for any couple-based IPV-adjacent \seqsplit{intervention} in this angle's search; the strongest RCT (Taft 2024) is a targeted secondary-prevention trial in a screened military population, not a universal premarital population.\\
\\[3pt]
Premarital-specific trials & decision-autonomy-measures & No study evaluating premarital or marriage-\seqsplit{preparation} education \seqsplit{specifically} through a decision-control or \seqsplit{reproductive}-autonomy outcome measure (rather than knowledge\slash{}attitude\slash{}\seqsplit{satisfaction}) was identified; the ARCHES trial (Miller 2016) is the closest analogue but is a clinic-based, not marriage-\seqsplit{preparation}, \seqsplit{intervention}.\\
\\[3pt]
Premarital-specific trials & decision-autonomy-measures & No study combining religious\slash{}faith-based marriage \seqsplit{preparation} \seqsplit{programming} with a validated decision-autonomy or coercion-\seqsplit{recognition} outcome measure was found in PubMed; such programmes, if evaluated at all, appear to be published outside PubMed-indexed venues or use non-\seqsplit{standardized} outcome measures.\\
\\[3pt]
Premarital-specific trials & digital-self-assessment-ipv & No PubMed-indexed record found \seqsplit{specifically} combining 'premarital' status (not-yet-married couples) with a scored anonymous self-check tool - all identified trials enrol women already \seqsplit{identifying} current\slash{}recent partner abuse, not a general premarital population screening before any disclosed problem exists.\\
\\[3pt]
Premarital-specific trials & ipv-screening-response & No PubMed-indexed evidence located specific to premarital\slash{}pre-nikah education settings combined with IPV screening\slash{}response protocols -- all direct LIVES\slash{}ARCHES \seqsplit{effectiveness} evidence (PMIDs 41267027, 41151831, 32460786) comes from antenatal\slash{}family-planning clinical encounters, not premarital or religious pre-marriage \seqsplit{counselling} contexts. This is a genuine evidence gap for T-PHE's specific setting, not merely a s...\\
\\[3pt]
Premarital-specific trials & islamic-medical-ethics-tier-one & No PubMed-indexed FIMA (Fiqh Council\slash{}Federation of Islamic Medical \seqsplit{Associations}) or IMANA formal position\slash{}guideline paper on premarital counseling, marriage \seqsplit{preparation}, or \seqsplit{reproductive}-decision consent was found (Q7 'FIMA\slash{}IMANA + guideline' returned 0 hits); the two Journal of IMA articles included (PMID 23610513, 23864762) are topical essays, not formal guidelines.\\
\\[3pt]
Premarital-specific trials & \seqsplit{manipulation}-literacy-\seqsplit{inoculation} & No PubMed-indexed record found that directly measures '\seqsplit{recognition} of coercive control\slash{}\seqsplit{gaslighting}' via a validated \seqsplit{psychometric} instrument \seqsplit{specifically} among \seqsplit{adolescents}\slash{}young adults in a premarital-education context (searches for 'coercive control \seqsplit{recognition} scale' and '\seqsplit{gaslighting} \seqsplit{recognition} education' returned 0 or near-0 direct hits; adjacent constructs like IPV \seqsplit{internalized} stigma (PMID 39536684) and gasli...\\
\\[3pt]
Premarital-specific trials & \seqsplit{manipulation}-literacy-\seqsplit{inoculation} & No PubMed-indexed study was found that directly tests whether \seqsplit{recognising}\slash{}labelling a \seqsplit{relationship} pattern as 'abuse' causally predicts subsequent help-seeking in a premarital or pre-\seqsplit{relationship} (as opposed to already-victimized) population -- the help-seeking records found (37882155, 38336660, 26442989) all study people who already \seqsplit{experienced} IPV, not people using an anonymous early-signal self-check before har...\\
\\[3pt]
Premarital-specific trials & peer-online-support-digital-safe-space & No PubMed-indexed record evaluates a premarital-education-linked or nikah\slash{}marriage-\seqsplit{preparation} digital peer community anywhere; this proposal's Layer 3 model appears untested in the indexed literature.\\
\\[3pt]
Premarital-specific trials & peer-online-support-digital-safe-space & No LMIC\slash{}Southeast Asia-specific systematic review of moderated digital peer support for \seqsplit{relationship} or premarital health topics was found via this angle; the closest LMIC-adjacent records (PMID 37540713, 38558241, 39116286) concern general digital health promotion, not peer community moderation \seqsplit{specifically}, and were excluded as off-angle.\\
\\[3pt]
Premarital-specific trials & premarital-education-outcomes & No PubMed-indexed study was found that evaluates a premarital or marriage education programme against decision-control, coercion-\seqsplit{recognition}, or informed-choice outcomes as primary endpoints (the specific outcome domain the FAILURE CRITERION targets) -- every \seqsplit{effectiveness} study located, including the Hawkins\slash{}ACF meta-analyses, measures \seqsplit{satisfaction}, \seqsplit{communication} skills, \seqsplit{relationship} stability, or quality of life...\\
\\[3pt]
Premarital-specific trials & premarital-education-outcomes & No PubMed-indexed evaluation of a STOP RULE-type screen-and-defer mechanism for suspected coercion\slash{}violence within premarital (as opposed to couple\slash{}marital) education \seqsplit{specifically} was found; the strongest supporting evidence (Markman 2013 PMID 23421844; Williamson 2015 PMID 25664643) is indirect, showing that joint \seqsplit{relationship} education can worsen outcomes for high-aggression\slash{}high-risk couples, not a direct test ...\\
\\[3pt]
Premarital-specific trials & premarital-education-outcomes & No PubMed-indexed premarital-education outcome study directly measures 'referral timeliness, \seqsplit{accessibility}, and follow-through' (P5) as a composite outcome; Williamson 2019 (PMID 30489129) and Williamson 2014 (PMID 24294929) are the closest available evidence, addressing help-seeking barriers and gateway effects \seqsplit{respectively}, but neither evaluates an actual referral pathway or its quality.\\
\\[3pt]
Premarital-specific trials & premarital-education-outcomes & No meta-analysis or systematic review \seqsplit{specifically} on faith-based\slash{}religious-leader-delivered premarital education and safety\slash{}coercion outcomes was found via the systematic-review-filtered queries (query 2 returned 0 hits on the exact phrasing tried); this may reflect true absence of such reviews in PubMed rather than a search-term failure, but should not be treated as confirmed absence without a broader grey-liter...\\
\\[3pt]
Premarital-specific trials & primary-prevention-frameworks & No PubMed-indexed paper was found that evaluates a premarital or pre-nuptial education programme \seqsplit{specifically} structured around a governance\slash{}decision-control framework comparable to T-PHE; the closest analogues (SASA!, \seqsplit{Indashyikirwa}, MFP) are ongoing-couple or community-level \seqsplit{interventions}, not marriage-transition-timed \seqsplit{interventions}.\\
\\[3pt]
Premarital-specific trials & primary-prevention-frameworks & Cost-\seqsplit{effectiveness} and \seqsplit{implementation}-science evidence located pertains to community \seqsplit{mobilisation} and CHW-delivered parenting\slash{}couples programmes, not premarital education delivered through religious \seqsplit{institutions}, limiting direct \seqsplit{transferability} to T-PHE's delivery channel.\\
\\[3pt]
Decision-control \seqsplit{instruments} & couple-work-safety-coercive-control & No PubMed-indexed study directly tests a 'stop rule' construct (halt joint work pending \seqsplit{confidential} assessment, never resume if coercive control is confirmed) as an explicit \seqsplit{intervention} design feature; the closest analogues (McCollum 2008, Schneider 2014, Karakurt 2016) describe screening\slash{}safety-\seqsplit{modification} practice \seqsplit{recommendations} rather than a validated stop\slash{}resume protocol.\\
\\[3pt]
Decision-control \seqsplit{instruments} & couple-work-safety-coercive-control & Evidence on screening\slash{}assessing male \seqsplit{perpetration} (rather than \seqsplit{victimization}) is \seqsplit{comparatively} sparse (Portnoy 2020), which matters for T-PHE's \seqsplit{bidirectional}\slash{}decision-control framing (P1) where either partner could be at risk.\\
\\[3pt]
Decision-control \seqsplit{instruments} & decision-autonomy-measures & No PubMed-indexed instrument was found that measures 'decision control' \seqsplit{specifically} at the point of a marriage-transition decision (proceed\slash{}delay\slash{}refuse marriage) for adult couples; the closest analogues are \seqsplit{reproductive}-decision \seqsplit{instruments} (RAS, RCS, FP-ADM) or an adolescent \seqsplit{empowerment} subscale ('choice of partners, marriage, and children' in Upadhyay 2021), none validated for the marriage-entry decision itself.\\
\\[3pt]
Decision-control \seqsplit{instruments} & decision-autonomy-measures & No instrument \seqsplit{specifically} measuring 'coercive control \seqsplit{recognition}' as a discrete competency (as opposed to experience\slash{}\seqsplit{victimization} scales like the RCS) was found; the T-PHE STOP RULE presumes \seqsplit{practitioners} and\slash{}or \seqsplit{participants} can recognize coercive control, but a dedicated \seqsplit{recognition}-literacy instrument (as distinct from an experience-report scale) was not located in this search.\\
\\[3pt]
Decision-control \seqsplit{instruments} & digital-self-assessment-ipv & No record was found addressing a combined single self-check instrument spanning fear, control, pressure, \seqsplit{reproductive} health AND mental health \seqsplit{simultaneously} (as T-PHE's Layer 0 proposes) - the literature covers these domains via separate, purpose-specific tools (IPV decision aids vs Danger Assessment vs \seqsplit{reproductive}-coercion measures), not one integrated scored instrument.\\
\\[3pt]
Decision-control \seqsplit{instruments} & ipv-screening-response & Evidence on harms of screening (adverse effects, loss of privacy, distress) is present (PMID 22565034, 26200817) but almost entirely measured only \seqsplit{immediately} post-screening, with very limited longer-term harm \seqsplit{surveillance} across the literature -- relevant to T-PHE's dignity\slash{}consent safeguard but not fully answered by current evidence.\\
\\[3pt]
Layer 0 \slash{} Layer 2 evidence & digital-self-assessment-ipv & Safety-harm data (partner discovery, \seqsplit{retaliation}) are addressed \seqsplit{qualitatively}\slash{}\seqsplit{methodologically} (Tarzia 2017) but no record reports a quantified adverse-event rate from an anonymous online self-check \seqsplit{specifically}, which the T-PHE failure criterion would need.\\
\\[3pt]
Layer 0 \slash{} Layer 2 evidence & peer-online-support-digital-safe-space & Evidence on moderator training, competence, and error rates in routing acute-risk \seqsplit{disclosures} (self-harm, IPV, abuse) to services is sparse (PMID 34193497 found only 5 usable studies after screening 397); no record quantifies false-negative or delayed-routing rates for non-clinical moderators.\\
\\[3pt]
Layer 0 \slash{} Layer 2 evidence & peer-online-support-digital-safe-space & Evidence \seqsplit{quantifying} how often anonymity itself is defeated by partner technology-\seqsplit{facilitated} \seqsplit{surveillance} of a specific digital peer-support platform (as opposed to general social media\slash{}apps) was not found; existing TFA evidence (PMID 35537445, 32998673, 40271872) describes general digital abuse mechanisms, not platform-specific breach of anonymous peer-support spaces.\\
\\[3pt]
\end{longtable}


\section{Discussion}

\subsection{What this map does and does not show}
This map shows that a structured, logged search across 18 categories, verified by an independent
adversarial pass, assembled a library of 313 unique PubMed-indexed records plus 17 named
guidance documents (330 unique records in total), of which 109 unique records are linked to at
least one of \tphe{}'s seven
propositions. It shows where that literature is comparatively dense (P1, P5) and where it is thin
(P6, P7), and that no record in the library is tagged as directly contradicting any proposition
--- itself notable, since it means every proposition currently has a body of supporting,
complicating, or merely relevant literature, but none has been tested directly against a
competing intervention. It does not show, and is not designed to show, that \tphe{} or any of its
five rules is effective, safe, or superior to an alternative: no admitted record evaluates \tphe{}
itself, and no record identified in this search evaluates any premarital-education programme, of
any design, in a Thai or Thai Muslim population.

\subsection{The claim ceiling}
Nothing in this document licenses an effectiveness inference. The library was built to ground a
programme theory in existing, adjacent, and analogical literature, and to make explicit, by
proposition, what remains untested. Where a table above shows a proposition with many
``supports''-tagged records, that density reflects the volume of adjacent literature the search
strategy located, not a strength-of-evidence judgement for \tphe{} itself; the corresponding
finding cells, read individually, make clear that almost none of that literature was generated on
a premarital, Thai, or Thai Muslim population, and none evaluates the governance layer \tphe{}
proposes.

\subsection{How this map will be used}
The gaps charted in \S3.4 --- absence of Thai/Thai-Muslim premarital-specific trials, absence of a
validated decision-control instrument for the marriage-transition decision, and absence of any
outcome evidence for the two proposed extension layers --- are exactly the gaps a future full
systematic review, or a primary study, would need to either fill or work around. The intended next
use of this map is to pre-specify the scope, eligibility criteria, and outcome set of a future full
systematic review of one or more of P1--P7 (most plausibly P1 or P5, where the adjacent evidence
base is largest), with a registered protocol (PROSPERO), independent duplicate screening, and a
risk-of-bias assessment (ROBIS for included systematic reviews, AMSTAR-2 for included reviews more
generally, Cochrane RoB~2 for included randomised trials) applied at the individual-study level ---
none of which this document performed or claims to have performed.

\section*{Declarations}
\addcontentsline{toc}{section}{Declarations}

\noindent\textbf{Conflict of interest.} The author is the founder of ARAYA Nikah Social Enterprise
Co., Ltd., the organisation whose premarital education programme motivated the \tphe{} programme
theory this evidence map supports. This is a direct conflict of interest, disclosed here as it is
in the \tphe{} full paper.

\noindent\textbf{Use of AI.} AI tools assisted literature search execution (NCBI E-utilities
queries), abstract screening, drafting of findings and relevance annotations, independent
verification re-fetching, and manuscript drafting, including generation of the tables in
\S3 directly from the logged JSON files. The author verified the generated counts against the
underlying files, ran the accompanying scripts, and is accountable for the content of this
document. No participant data were used at any stage; every source is a published or
publicly-issued document (a journal article, its PMID, or a named guidance body's publication).

\noindent\textbf{Data availability.} The complete evidence library --- every search log, every
verification log, the admitted-entry table, the knowledge graph --- is public in the repository
(\texttt{library/}, https://github.com/morrocwi/RE\_T-PHE). The two scripts that produce every count
and table in this document, and this manuscript's source, are released in the same repository under
\texttt{paper/evidence-map/} after the independent adversarial review that the repository requires
before any public change.

\noindent\textbf{Funding.} No external funding was received for this evidence map.

\begin{thebibliography}{999}\footnotesize
% Auto-generated by generate_tex_fragments.py -- do not hand-edit.
% \bibitem text copied verbatim from paper/sections/refs_all.tex
% for every reference key this document cites.
\bibitem{ref1} GBD 2023 IPV and SVAC Collaborators. Disease burden attributable to intimate partner violence against females and sexual violence against children in 204 countries and territories, 1990-2023: a systematic analysis for the Global Burden of Disease Study 2023. Lancet (London, England). 2026. PMID: 41386261.
\bibitem{ref2} Sardinha L. Global, regional, and national prevalence estimates of physical or sexual, or both, intimate partner violence against women in 2018. Lancet (London, England). 2022. PMID: 35182472.
\bibitem{ref3} Kuswanto H. Prevalence of and factors associated with female child marriage in Indonesia. PLoS ONE. 2024. PMID: 38968277.
\bibitem{ref4} Rumble L. An empirical exploration of female child marriage determinants in Indonesia. BMC Public Health. 2018. PMID: 29587705.
\bibitem{ref5} Islam MK. Attitudes to and experiences of intimate partner violence among Rohingya women who married before eighteen years of age. Global Health Action. 2021. PMID: 34323166.
\bibitem{ref6} Hawkins AJ. Does marriage and relationship education work? A meta-analytic study. Journal of Consulting and Clinical Psychology. 2008. PMID: 18837590.
\bibitem{ref9} Markman HJ. A randomized clinical trial of the effectiveness of premarital intervention: moderators of divorce outcomes. Journal of Family Psychology. 2013. PMID: 23421844.
\bibitem{ref10} Williamson HC. Risk moderates the outcome of relationship education: a randomized controlled trial. Journal of Consulting and Clinical Psychology. 2015. PMID: 25664643.
\bibitem{ref16} Upadhyay Ushma D. Development and validation of a reproductive autonomy scale. Studies in Family Planning. 2014. PMID: 24615573.
\bibitem{ref17} McCauley Heather L. Psychometric properties and refinement of the Reproductive Coercion Scale. Contraception. 2017. PMID: 27639927.
\bibitem{ref18} Upadhyay Ushma D. Development and Validation of the Sexual and Reproductive Empowerment Scale for Adolescents and Young Adults. The Journal of Adolescent Health. 2021. PMID: 32690468.
\bibitem{ref19} Dagher Myriam. Adaptation and psychometric assessment of a sexual and reproductive empowerment scale in Arabic among refugee and non-refugee adolescent girls. BMC Medical Research Methodology. 2024. PMID: 39266993.
\bibitem{ref20} Alomair N. Factors influencing sexual and reproductive health of Muslim women: a systematic review. Reproductive Health. 2020. PMID: 32138744.
\bibitem{ref24} Abramsky T. Findings from the SASA! Study: a cluster randomized controlled trial to assess the impact of a community mobilization intervention to prevent violence against women and reduce HIV risk in Kampala, Uganda. BMC medicine. 2014. PMID: 25248996.
\bibitem{ref27} Shaw B. Shifting Norms in Faith Communities to Reduce Intimate Partner Violence: Results from a Cluster Randomized Controlled Trial in Nigeria. Journal of interpersonal violence. 2023. PMID: 37329160.
\bibitem{ref33} Hawkins AJ. Exploring programmatic moderators of the effectiveness of marriage and relationship education programs: a meta-analytic study. Behavior Therapy. 2012. PMID: 22304880.
\bibitem{ref38} O'Doherty L. Screening women for intimate partner violence in healthcare settings. Cochrane Database of Systematic Reviews. 2015. PMID: 26200817.
\bibitem{ref40} US Preventive Services Task Force. Screening for Intimate Partner Violence and Caregiver Abuse of Older or Vulnerable Adults: US Preventive Services Task Force Recommendation Statement. JAMA. 2025. PMID: 40553450.
\bibitem{ref42} Betron. Effectiveness of a clinic-based counselling intervention on risk of experiencing intimate partner violence and reproductive coercion: a matched-pair cluster-controlled trial in Ebonyi and Sokoto, Nigeria. BMJ Global Health. 2025. PMID: 41151831.
\bibitem{ref45} Karakurt. Couples Therapy for Intimate Partner Violence: A Systematic Review and Meta-Analysis. Journal of marital and family therapy. 2016. PMID: 27377617.
\bibitem{ref46} McCollum. Couples treatment for interpersonal violence: a review of outcome research literature and current clinical practices. Violence and victims. 2008. PMID: 18624089.
\bibitem{ref47} Schneider. From scared to repaired: using an attachment-based perspective to understand situational couple violence. Journal of marital and family therapy. 2014. PMID: 25039387.
\bibitem{ref49} Robertson. Women and men's use of coercive control in intimate partner violence. Violence and victims. 2011. PMID: 21780535.
\bibitem{ref51} Evans. "Even 'Daily' is Not Enough": How Well Do We Measure Domestic Violence and Abuse?-A Think-Aloud Study of a Commonly Used Self-Report Scale. Violence and victims. 2016. PMID: 26645540.
\bibitem{ref92} Taft. Examining strength at home couples to prevent intimate partner violence on a military installation: A randomized controlled trial. Journal of consulting and clinical psychology. 2024. PMID: 38206858.
\bibitem{ref93} Ghuman SJ. Women's autonomy and child survival: a comparison of Muslims and non-Muslims in four Asian countries. Demography. 2003. PMID: 12962056.
\bibitem{ref94} Chen. Intimate partner violence: counseling, community resources, and legal issues for IPV victims and perpetrators. FP essentials. 2013. PMID: 24053261.
\bibitem{ref95} Decker. Implementing Trauma-Informed Partner Violence Assessment in Family Planning Clinics. Journal of women's health. 2017. PMID: 28375750.
\bibitem{ref96} MacMillan. Approaches to screening for intimate partner violence in health care settings: a randomized trial. JAMA. 2006. PMID: 16882959.
\bibitem{ref97} Ford-Gilboe. Development of a brief measure of intimate partner violence experiences: the Composite Abuse Scale (Revised)-Short Form (CASR-SF). BMJ open. 2016. PMID: 27927659.
\bibitem{ref98} Hegarty. Validation of the Coercive-Composite Abuse Scale (C-CAS) in an Australian sample of women experiencing intimate partner abuse. BMJ open. 2026. PMID: 42323273.
\bibitem{ref99} Wilson. Psychometric Properties of the Coercion in Intimate Partner Relationships Scale. Assessment. 2023. PMID: 34167331.
\bibitem{ref100} Sarac A. Cultural adaptation and pilot evaluation of the Couple Learning Program (CLP) among young Turkish couples: A mixed-method study. Acta Psychologica. 2026. PMID: 42442188.
\bibitem{ref101} Welland C. Culturally specific treatment for partner-abusive Latino men: a qualitative study to identify and implement program components. Violence and Victims. 2010. PMID: 21287968.
\bibitem{ref102} Sikoki B. A qualitative study on perceptions of adolescents' sexual and reproductive health education in Yogyakarta, Indonesia. International Journal of Adolescent Medicine and Health. 2024. PMID: 39395200.
\bibitem{ref103} Fauk NK. Cultural and religious determinants of HIV transmission: A qualitative study with people living with HIV in Belu and Yogyakarta, Indonesia. PLoS ONE. 2021. PMID: 34780506.
\bibitem{ref104} Grisurapong S. Establishing a one-stop crisis center for women suffering violence in Khonkaen hospital, Thailand. International Journal of Gynaecology and Obstetrics. 2002. PMID: 12429436.
\bibitem{ref105} Nagae M. [Domestic violence as a women's health issue--activities of health professionals in Thailand]. Japanese Journal of Public Health (Nihon Koshu Eisei Zasshi). 2004. PMID: 15162975.
\bibitem{ref106} Williamson HC. Does premarital education decrease or increase couples' later help-seeking? Journal of Family Psychology. 2014. PMID: 24294929.
\bibitem{ref107} Williamson HC. Barriers and facilitators of relationship help-seeking among low-income couples. Journal of Family Psychology. 2019. PMID: 30489129.
\bibitem{ref108} Oduenyi. Integrating gender-based violence first-line response into family planning and antenatal care services: an assessment of feasibility, appropriateness and acceptability to healthcare providers in Sokoto and Ebonyi States, Nigeria. BMC Health Services Research. 2025. PMID: 41267027.
\bibitem{ref109} Rhodes. The anatomy of a community health center system-level intervention for intimate partner violence. Journal of Urban Health. 2014. PMID: 23917943.
\bibitem{ref110} Miller Elizabeth. A family planning clinic-based intervention to address reproductive coercion: a cluster randomized controlled trial. Contraception. 2016. PMID: 26892333.
\bibitem{ref111} Zaher. Effect of domestic violence training: systematic review of randomized controlled trials. Canadian Family Physician. 2014. PMID: 25022633.
\bibitem{ref112} Jacubowski N. Marriage is not a safe place: heterosexual marriage and HIV-related vulnerability in Indonesia. Culture, Health \& Sexuality. 2008. PMID: 18038283.
\bibitem{ref113} Michau L. SASA! Together: An evolution of the SASA! approach to prevent violence against women. Evaluation and program planning. 2021. PMID: 33578229.
\bibitem{ref114} Stern E. Lessons learned from implementing Indashyikirwa in Rwanda - an adaptation of the SASA! approach to prevent and respond to intimate partner violence. Evaluation and program planning. 2018. PMID: 30125773.
\bibitem{ref115} Clyde TL. Revising Premarital Relationship Interventions for the Next Generation. Journal of Marital and Family Therapy. 2020. PMID: 30725473.
\bibitem{ref116} Silliman B. Improving practice in marriage preparation. Journal of Sex \& Marital Therapy. 1999. PMID: 10081741.
\bibitem{ref118} Zust BL. 10-Year Study of Christian Church Support for Domestic Violence Victims: 2005-2015. Journal of Interpersonal Violence. 2021. PMID: 33729071.
\bibitem{ref160} Schenck D et al. Ethical Analysis and Policy Recommendations Regarding Domino Liver Transplantation. Transplantation. 2018; PMID: 29708521.
\bibitem{ref164} Messing JT et al. The average predictive validity of intimate partner violence risk assessment instruments. Journal of Interpersonal Violence. 2013; PMID: 23262817.
\bibitem{ref165} Wang PL et al. Assessing the Danger: Validation of Taiwan Intimate Partner Violence Danger Assessment. Journal of Interpersonal Violence. 2015; PMID: 25315482.
\bibitem{ref168} Gerth J et al. [The Ontario Domestic Assault Risk Assessment (ODARA) - Validity und authorised German translation of an intimate partner violence screening tool]. Fortschritte der Neurologie-Psychiatrie. 2014; PMID: 25383928.
\bibitem{ref171} Hogan NR et al. Investigating racial disparities in violence risk assessment using the Spousal Assault Risk Assessment Guide-Version 3. Psychological Assessment. 2024; PMID: 38512165.

% Auto-generated by generate_tex_fragments.py -- do not hand-edit.
% Vancouver-style bibitems for library records cited in this document that
% are not among the \bibitem entries copied verbatim from
% paper/sections/refs_all.tex. Built directly from library/library.json
% fields (title, journal, year, pmid); the author field is taken from a cached
% PubMed esummary record (pubmed_authors_cache.json) keyed by PMID, first author
% 'Surname XY' + ' et al.' when more than one author, institutional/group authors
% kept verbatim, falling back to the library.json first_author field only when no
% cached esummary record is available for that PMID.
\bibitem{srp14996680} U. S. Preventive Services Task Force. Screening for family and intimate partner violence: \seqsplit{recommendation} statement. Annals of Internal Medicine. 2004. PMID: 14996680.
\bibitem{srp17457732} Blackwood E. Regulation of sexuality in Indonesian discourse: normative gender, criminal law and shifting strategies of control. Culture, Health \& Sexuality. 2007. PMID: 17457732.
\bibitem{srp21355646} Owen JJ et al. The role of leaders' working alliance in premarital education. Journal of Family Psychology. 2011. PMID: 21355646.
\bibitem{srp22565034} Nelson HD et al. Screening women for intimate partner violence: a systematic review to update the U.S. Preventive Services Task Force \seqsplit{recommendation}. Annals of Internal Medicine. 2012. PMID: 22565034.
\bibitem{srp24860075} Tiwari A et al. Evaluating the Chinese Revised \seqsplit{Controlling} Behaviors Scale. Journal of \seqsplit{interpersonal} violence. 2015. PMID: 24860075.
\bibitem{srp25467583} García-Moreno C et al. The health-systems response to violence against women. Lancet. 2015. PMID: 25467583.
\bibitem{srp26924488} Michaels-Igbokwe C et al. Cost and cost-\seqsplit{effectiveness} analysis of a community \seqsplit{mobilisation} \seqsplit{intervention} to reduce intimate partner violence in Kampala, Uganda. BMC public health. 2016. PMID: 26924488.
\bibitem{srp26982668} Bhandari TR et al. \seqsplit{Construction} and Validation of a Women's Autonomy \seqsplit{Measurement} Scale with Reference to \seqsplit{Utilization} of Maternal Health Care Services in Nepal. JNMA; Journal of the Nepal Medical \seqsplit{Association}. 2014. PMID: 26982668.
\bibitem{srp27682273} Starmann E et al. Exploring Couples' Processes of Change in the Context of SASA!, a Violence Against Women and HIV Prevention \seqsplit{Intervention} in Uganda. Prevention science. 2017. PMID: 27682273.
\bibitem{srp28649297} Sprague S et al. A scoping review of intimate partner violence assistance programmes within health care settings. European journal of \seqsplit{psychotraumatology}. 2017. PMID: 28649297.
\bibitem{srp28789659} Thombs BD et al. \seqsplit{Consistency} and sources of divergence in \seqsplit{recommendations} on screening with \seqsplit{questionnaires} for presently \seqsplit{experienced} health problems or symptoms: a comparison of \seqsplit{recommendations} from the Canadian Task Force on Preventive Health Care, UK National Screening Committee, and US Preventive Services Task Force. BMC Medicine. 2017. PMID: 28789659.
\bibitem{srp29162117} Semahegn A et al. Community based \seqsplit{intervention} to prevent domestic violence against women in the \seqsplit{reproductive} age in \seqsplit{Northwestern} Ethiopia: a study protocol for community based cluster randomized trial. \seqsplit{Reproductive} health. 2017. PMID: 29162117.
\bibitem{srp29384426} Choi YJ et al. Korean clergy for healthy families: online \seqsplit{intervention} for preventing intimate partner violence. Global health promotion. 2019. PMID: 29384426.
\bibitem{srp29683049} Bundock K et al. \seqsplit{Adolescents}' Help-Seeking Behavior and Intentions Following Adolescent Dating Violence: A Systematic Review. Trauma, Violence \& Abuse. 2020. PMID: 29683049.
\bibitem{srp30357304} Feltner C et al. Screening for Intimate Partner Violence, Elder Abuse, and Abuse of Vulnerable Adults: Evidence Report and Systematic Review for the US Preventive Services Task Force. JAMA. 2018. PMID: 30357304.
\bibitem{srp30357305} US Preventive Services Task Force. Screening for Intimate Partner Violence, Elder Abuse, and Abuse of Vulnerable Adults: US Preventive Services Task Force Final \seqsplit{Recommendation} Statement. JAMA. 2018. PMID: 30357305.
\bibitem{srp30613427} Abramsky T et al. Changing the norms that drive intimate partner violence: findings from a cluster randomised trial on what \seqsplit{predisposes} bystanders to intervene. BMJ global health. 2018. PMID: 30613427.
\bibitem{srp31012540} Lee ASD et al. Improving Provider Readiness for Intimate Partner Violence Screening. Worldviews on evidence-based nursing. 2019. PMID: 31012540.
\bibitem{srp31306331} Hill AL et al. \seqsplit{Reproductive} Coercion and \seqsplit{Relationship} Abuse Among \seqsplit{Adolescents} and Young Women Seeking Care at School Health Centers. Obstetrics and Gynecology. 2019. PMID: 31306331.
\bibitem{srp31686618} McLean L et al. Shifting and \seqsplit{transforming} gender-\seqsplit{inequitable} beliefs, behaviours and norms in intimate \seqsplit{partnerships}: the \seqsplit{Indashyikirwa} couples' programme in Rwanda. Culture, health \& sexuality. 2020. PMID: 31686618.
\bibitem{srp31896308} Couture-Carron A. Shame, Family Honor, and Dating Abuse: Lessons From an \seqsplit{Exploratory} Study of South Asian Muslims. Violence Against Women. 2020. PMID: 31896308.
\bibitem{srp32460786} Uysal J et al. Protocol for a matched-pair cluster control trial of ARCHES (Addressing \seqsplit{Reproductive} Coercion in Health Settings) among women and girls seeking \seqsplit{contraceptive} services from community-based clinics in Nairobi, Kenya. \seqsplit{Reproductive} Health. 2020. PMID: 32460786.
\bibitem{srp32608505} Hameed M et al. \seqsplit{Psychological} therapies for women who experience intimate partner violence. Cochrane Database of Systematic Reviews. 2020. PMID: 32608505.
\bibitem{srp32791967} Portnoy GA et al. Patient and provider barriers, \seqsplit{facilitators}, and \seqsplit{implementation} \seqsplit{preferences} of intimate partner violence \seqsplit{perpetration} screening. BMC health services research. 2020. PMID: 32791967.
\bibitem{srp32869478} Stangl AL et al. Is HIV index testing and partner \seqsplit{notification} safe for adolescent girls and young women in low- and middle-income countries? Journal of the \seqsplit{International} AIDS Society. 2020. PMID: 32869478.
\bibitem{srp32896204} Stern E et al. Pathways of change: \seqsplit{qualitative} \seqsplit{evaluations} of intimate partner violence prevention programmes in Ghana, Rwanda, South Africa and Tajikistan. Culture, health \& sexuality. 2021. PMID: 32896204.
\bibitem{srp33336831} Wood SN et al. \seqsplit{Reproductive} Coercion among Intimate Partner Violence Survivors in Nairobi. Studies in Family Planning. 2020. PMID: 33336831.
\bibitem{srp34253076} Alghamdi MS et al. Intimate Partner Violence among Canadian Muslim Women. Journal of \seqsplit{Interpersonal} Violence. 2022. PMID: 34253076.
\bibitem{srp34527969} Islam MM et al. Factors affecting child marriage and \seqsplit{contraceptive} use among Rohingya girls in refugee camps. The Lancet Regional Health - Western Pacific. 2021. PMID: 34527969.
\bibitem{srp34770188} James LE et al. \seqsplit{Development} and Testing of a Community-Based \seqsplit{Intervention} to Address Intimate Partner Violence among Rohingya and Syrian Refugees: A Social Norms-Based Mental Health-Integrated Approach. \seqsplit{International} Journal of \seqsplit{Environmental} Research and Public Health. 2021. PMID: 34770188.
\bibitem{srp34838218} Yakubovich AR et al. Housing \seqsplit{interventions} for women \seqsplit{experiencing} intimate partner violence: a systematic review. Lancet Public Health. 2022. PMID: 34838218.
\bibitem{srp35611345} Gutowski ER et al. Coercive Control in the Courtroom: the Legal Abuse Scale (LAS). Journal of family violence. 2023. PMID: 35611345.
\bibitem{srp36163286} \seqsplit{Ratnaningsih} M et al. Child Marriage \seqsplit{Acceptability} Index (CMAI) as an essential indicator: an \seqsplit{investigation} in South and Central Sulawesi, Indonesia. Global Health Research and Policy. 2022. PMID: 36163286.
\bibitem{srp36429857} Subramanee SD et al. Child Marriage in South Asia: A Systematic Review. \seqsplit{International} Journal of \seqsplit{Environmental} Research and Public Health. 2022. PMID: 36429857.
\bibitem{srp36550423} Gausman J et al. Girl child marriage and the social context of \seqsplit{displacement}: a \seqsplit{qualitative} \seqsplit{comparative} \seqsplit{exploration} of Syrian refugees in Jordan and Rohingya refugees in Bangladesh. BMC Public Health. 2022. PMID: 36550423.
\bibitem{srp36657958} Riches E et al. Evaluation of the \seqsplit{psychometric} properties of the \seqsplit{Reproductive} Autonomy Scale for use in the UK. BMJ Sexual \& \seqsplit{Reproductive} Health. 2023. PMID: 36657958.
\bibitem{srp36693203} Singh V et al. Family medicine provider and staff \seqsplit{identification} and response to male patient intimate partner violence \seqsplit{perpetration}. Annals of Family Medicine. 2022. PMID: 36693203.
\bibitem{srp37316890} Boyce SC et al. Measuring social norms of intimate partner violence to exert control over wife agency, sexuality, and \seqsplit{reproductive} autonomy: an item response modelling of the IPV-ASRA scale. \seqsplit{Reproductive} Health. 2023. PMID: 37316890.
\bibitem{srp37649394} Debinski B et al. \seqsplit{Qualitative} Assessment of Key \seqsplit{Implementation} Factors in a Faith-Based Response to Intimate Partner Violence. Health promotion practice. 2025. PMID: 37649394.
\bibitem{srp37922257} Gausman J et al. Validation of a measure to assess decision-making autonomy in family planning services in three low- and middle-income countries: The Family Planning Autonomous Decision-Making scale (FP-ADM). PLOS ONE. 2023. PMID: 37922257.
\bibitem{srp38339841} Hazar S et al. Turkish adaptation of the \seqsplit{reproductive} autonomy scale: Validity and \seqsplit{reliability} study. The Journal of Obstetrics and \seqsplit{Gynaecology} Research. 2024. PMID: 38339841.
\bibitem{srp38970860} Han X et al. Maternal \seqsplit{victimization} and neglected offspring: Child marriage, IPV and depression symptoms among Salar Muslim women. Child Abuse \& Neglect. 2024. PMID: 38970860.
\bibitem{srp39425490} Sharifnia AM et al. Muslim Women's \seqsplit{Experiences} of Domestic Violence and Abuse: A Meta-\seqsplit{Ethnography} of Global Evidence. Trauma, Violence \& Abuse. 2025. PMID: 39425490.
\bibitem{srp39757299} Park J et al. \seqsplit{Videoconferencing} Delivery of the Seoul Premarital Education Program During COVID-19: A Quasi-\seqsplit{experimental} Study Using Inverse \seqsplit{Probability} of Treatment Weighting. Prevention Science. 2026. PMID: 39757299.
\bibitem{srp40091116} Mieth K et al. "What other option did I have?"- The effect of conflict and \seqsplit{displacement} on child marriage and early \seqsplit{childbearing} among displaced Rohingya \seqsplit{adolescents}. Conflict and Health. 2025. PMID: 40091116.
\bibitem{srp40116416} Azalee M et al. Wounded Healers: Exploring Coping With Intimate Partner Violence Among Health Care Workers in Malaysia. Asia-Pacific Journal of Public Health. 2025. PMID: 40116416.
\bibitem{srp41330529} Giles F et al. Women's \seqsplit{preferences} for how health \seqsplit{practitioners} respond to coercive control by a partner: Open-ended survey \seqsplit{qualitative} analysis. Australian Journal of General Practice. 2025. PMID: 41330529.
\bibitem{srp41454335} \seqsplit{Nsengiyumva} R et al. Perceived need of \seqsplit{information} and education about conception \seqsplit{preparedness} (CP) among engaged and soon to be married couples: a \seqsplit{qualitative} survey. \seqsplit{Reproductive} Health. 2025. PMID: 41454335.
\bibitem{srp41613074} Böttcher B et al. The \seqsplit{Reproductive} Autonomy Scale: German validation and \seqsplit{application} in an Austrian student population. Frontiers in Public Health. 2025. PMID: 41613074.
\bibitem{srp42015126} Kurnaz K et al. \seqsplit{Reproductive} coercion, gender role attitudes and fertility desire among married women in Türkiye: a \seqsplit{descriptive} and \seqsplit{explanatory} study. BMC Public Health. 2026. PMID: 42015126.
\bibitem{srp42300599} Phonkacha N et al. Technology Access and Attitudes Toward Intimate Partner Violence: Evidence From a National Survey in Thailand. Violence Against Women. 2026. PMID: 42300599.
\bibitem{srp9779330} Sayers SL et al. Prevention of marital \seqsplit{dysfunction}: behavioral approaches and beyond. Clinical Psychology Review. 1998. PMID: 9779330.

\end{thebibliography}

\end{document}
